\documentclass[bachelor, fontsize=14]{diploma}

% ──────────────────────────────────────────────────────────────────────────
%  Дополнительные пакеты (в шаблоне diploma их нет)
% ──────────────────────────────────────────────────────────────────────────
\usepackage{tikz}
\usetikzlibrary{arrows.meta, positioning, calc, fit, backgrounds, shapes.geometric, shapes.multipart}
\usepackage{booktabs}
\usepackage{multirow}
\usepackage{siunitx}

% ──────────────────────────────────────────────────────────────────────────
%  TODO: подтвердить данные титульного листа у автора
% ──────────────────────────────────────────────────────────────────────────
\student{Кочкожаров Иван Вячеславович}      % TODO: ФИО полностью
\group{М8О-408Б-22}                           % TODO: номер учебной группы
\theme{Модуль интеграции колоночного файлового формата Vortex в открытый табличный формат и сравнение со стандартными файловыми форматами}

\supervisor{Гаврилов Евгений Сергеевич}                   % TODO: ФИО научного руководителя
\firstConsultant{}                            % TODO: консультант (если есть)
\reviewer{}                                   % TODO: ФИО рецензента

\faculty{№ 8 <<Компьютерные науки и прикладная математика>>}
\department{806}                              % TODO: номер кафедры
\speciality{01.03.02 <<Прикладная математика и информатика>>}  % TODO
\profile{Информатика}                         % TODO: профиль

\departmentFullName{№ 806}                    % TODO
\headOfDepartment{Крылов Сергей Сергеевич}                           % TODO: зав. кафедрой

\date{\uline{\hspace{24pt}} \the\year\ года}

% Перечень сокращений и обозначений. Порядок — по алфавиту.

\newacronym{alp}{ALP}{Adaptive Lossless floating-Point compression}
\newacronym{api}{API}{Application Programming Interface}
\newacronym{cbo}{CBO}{Cost-Based Optimizer}
\newacronym{cdc}{CDC}{Change Data Capture}
\newacronym{cpu}{CPU}{Central Processing Unit}
\newacronym{dwh}{DWH}{Data Warehouse}
\newacronym{fsst}{FSST}{Fast Static Symbol Table}
\newacronym{gpu}{GPU}{Graphics Processing Unit}
\newacronym{hdfs}{HDFS}{Hadoop Distributed File System}
\newacronym{jni}{JNI}{Java Native Interface}
\newacronym{jvm}{JVM}{Java Virtual Machine}
\newacronym{lsm}{LSM}{Log-Structured Merge-tree}
\newacronym{nrt}{NRT}{Near Real-Time}
\newacronym{olap}{OLAP}{Online Analytical Processing}
\newacronym{oltp}{OLTP}{Online Transaction Processing}
\newacronym{orc}{ORC}{Optimized Row Columnar}
\newacronym{rle}{RLE}{Run-Length Encoding}
\newacronym{s3}{S3}{Simple Storage Service}
\newacronym{simd}{SIMD}{Single Instruction, Multiple Data}
\newacronym{sql}{SQL}{Structured Query Language}
\newacronym{tpcds}{TPC-DS}{Transaction Processing Performance Council Decision Support benchmark}
\newacronym{vkr}{ВКР}{выпускная квалификационная работа}

% Термины и определения. Порядок — по алфавиту (кириллица, затем латиница).
% Нумерация id произвольная, но уникальная.

\newglossaryentry{t-encoding}{
    name={кодирование колонки},
    description={преобразование последовательности однотипных значений столбца в более компактное двоичное представление с использованием статистических свойств данных (повторов, ограниченного диапазона, общих префиксов и т. п.)}
}
\newglossaryentry{t-skipping}{
    name={пропуск данных},
    description={(англ. data skipping) приём оптимизации чтения, при котором по статистикам (минимум, максимум, число null-значений) блоков данных заранее определяется, что блок не содержит подходящих под предикат строк, и его чтение пропускается}
}
\newglossaryentry{t-compaction}{
    name={слияние (компакция)},
    description={(англ. compaction) фоновая операция в LSM-дереве, объединяющая несколько отсортированных файлов одного уровня в меньшее число файлов следующего уровня с удалением перезаписанных и удалённых записей}
}
\newglossaryentry{t-format}{
    name={файловый формат},
    description={спецификация двоичного представления табличных данных в файле, включающая разбиение на блоки, схему хранения значений (построчное или поколоночное), способы кодирования, сжатия и встроенные метаданные (статистики, индексы)}
}
\newglossaryentry{t-tableformat}{
    name={открытый табличный формат},
    description={(англ. open table format) спецификация метаданных поверх набора файлов данных, обеспечивающая представление их как единой таблицы с поддержкой транзакций, эволюции схемы, моментальных снимков и параллельного доступа; примеры — Apache Iceberg, Delta Lake, Apache Hudi, Apache Paimon}
}
\newglossaryentry{t-columnar}{
    name={колоночный формат хранения},
    description={формат, в котором значения одного столбца таблицы располагаются в файле непрерывно; обеспечивает высокую степень сжатия однотипных данных и чтение только запрашиваемых столбцов}
}
\newglossaryentry{t-predicate}{
    name={предикат},
    description={логическое условие фильтрации строк (например, \texttt{col > 100 AND col < 200}); может передаваться в файловый формат для отсечения блоков данных на уровне декодера (predicate pushdown)}
}
\newglossaryentry{t-mor}{
    name={слияние при чтении},
    description={(англ. merge-on-read) стратегия таблиц с первичным ключом, при которой обновления и удаления записываются как новые версии, а актуальное состояние строки вычисляется во время чтения объединением версий из всех уровней LSM-дерева по полю последовательности}
}
\newglossaryentry{t-lakehouse}{
    name={lakehouse},
    description={архитектура хранения данных, объединяющая дешёвое масштабируемое хранилище объектов озера данных (data lake) с транзакционными гарантиями, управлением схемой и SQL-доступом, характерными для хранилищ данных (data warehouse)}
}
\newglossaryentry{t-lsm}{
    name={LSM-дерево},
    description={(англ. Log-Structured Merge-tree) структура данных, оптимизированная под интенсивную запись: новые данные буферизуются в памяти и сбрасываются на диск отсортированными неизменяемыми файлами (уровень~0), которые затем фоновым слиянием перемещаются на нижележащие уровни}
}
\newglossaryentry{t-olap}{
    name={OLAP},
    description={(англ. Online Analytical Processing) класс нагрузок, характеризующихся сложными аналитическими запросами с агрегацией над большими объёмами данных и редкими изменениями; противопоставляется OLTP}
}
\newglossaryentry{t-oltp}{
    name={OLTP},
    description={(англ. Online Transaction Processing) класс нагрузок, характеризующихся большим числом коротких транзакций чтения и записи отдельных записей по ключу}
}


\addbibresource{main.bib}

\begin{document}
    \maketitle

    % \includepdf[pages=-]{extra/task} % Задание на ВКР (вкладывается отдельно)
    \setcounter{page}{2}

    \abstract % Структурный элемент: РЕФЕРАТ. Строка со сведениями об объёме формируется автоматически.

\keywords{APACHE PAIMON, КОЛОНОЧНЫЙ ФОРМАТ, VORTEX, PARQUET, ORC, LAKEHOUSE, LSM-ДЕРЕВО, КОДИРОВАНИЕ ДАННЫХ, ПРОПУСК ДАННЫХ, НАГРУЗОЧНОЕ ТЕСТИРОВАНИЕ, APACHE FLINK}

Объект исследования --- колоночные файловые форматы хранения данных в составе
lakehouse-системы Apache Paimon: Apache Parquet, Apache ORC и Vortex. Предмет
исследования --- алгоритмы кодирования, сжатия и пропуска данных, реализуемые
этими форматами, и их влияние на производительность чтения и записи таблиц с
первичным ключом и таблиц только для добавления.

Цель работы --- определить область применимости формата Vortex в Apache Paimon,
разработать модуль его интеграции и провести сравнительный анализ форматов на
представительных нагрузках.

В ходе работы выполнен аналитический обзор эволюции систем хранения аналитических
данных --- от хранилищ данных к архитектуре lakehouse, --- открытых табличных
форматов и колоночных файловых форматов; разработан и интегрирован в Apache Paimon
модуль поддержки формата Vortex: реализация интерфейса \texttt{FileFormat},
конвертер предикатов, взаимодействие с нативной библиотекой через Java Native
Interface; проведено функциональное тестирование артефакта в контейнерах;
развёрнуто окружение нагрузочного тестирования на кластере из пяти узлов
средствами Ansible; проведены две серии экспериментов --- batch-аналитика на
наборе TPC-DS масштаба \SI{100}{\giga\byte} и near-real-time-сценарий с потоковой
записью и параллельными аналитическими запросами; собраны и статистически
обработаны метрики производительности.

Результаты: на batch-аналитике (таблицы только для добавления, крупные файлы)
Vortex доминирует над Parquet --- быстрее на \num{82} запросах из \num{103}, среднее
ускорение порядка \SI{25}{\percent}, пиковое --- до \SI{80}{\percent}; именно этот
сценарий --- основная область применимости Vortex. На потоковой
near-real-time-нагрузке (таблица с первичным ключом, слияние при чтении) Vortex
эквивалентен Parquet и ORC по пропускной способности записи; по типичной задержке
чтения (медиане) Vortex --- быстрейший формат в обеих фазах и на большинстве
запросов (девять из тринадцати как в фазе FINAL, так и в фазе DURING), особенно ---
на запросах с селективным предикатом, диапазонным фильтром по строке-префиксу и в
многошаговой аналитике. На простой группировке и многомерной агрегации без
селективного предиката Vortex незначительно (в пределах единиц процентов) уступает
конкурентам --- следствие отсутствия проталкивания агрегатов в интерфейсе форматов
Apache Paimon. По хвостам задержки в фазе DURING ORC и Vortex показывают узкие
распределения (среднее почти равно медиане), Parquet --- более широкое
распределение с эпизодическими выбросами в десятки секунд в моменты сброса данных
по чекпойнту, что согласуется с архитектурным выигрышем Vortex за счёт независимого
нативного пула соединений с объектным хранилищем. Сформулированы рекомендации по
выбору формата хранения для различных классов нагрузок.

Область применения результатов --- проектирование витрин данных на базе Apache
Paimon, настройка lakehouse-хранилищ, развитие экосистемы Apache Paimon.
        % РЕФЕРАТ

    \tableofcontents                     % СОДЕРЖАНИЕ
    \termsanddefenitions                 % ТЕРМИНЫ И ОПРЕДЕЛЕНИЯ
    \listofabbreviations                 % ПЕРЕЧЕНЬ СОКРАЩЕНИЙ И ОБОЗНАЧЕНИЙ

    \introduction % Структурный элемент: ВВЕДЕНИЕ

Современные организации придерживаются стратегии накопления всех доступных данных в
расчёте на их будущее использование в задачах, которые на момент сбора ещё не
сформулированы. Эта практика породила потребность в дешёвом, масштабируемом и при
этом аналитически пригодном хранилище. Классические хранилища данных (data
warehouse) обеспечивают транзакционность, управление схемой и быстрый SQL-доступ, но
плохо масштабируются по стоимости и не приспособлены к слабоструктурированным
данным. Озёра данных (data lake) на основе распределённых файловых систем и
объектных хранилищ дёшевы и масштабируемы, но не дают транзакционных гарантий и
управления метаданными. Архитектура lakehouse~\cite{armbrust2021lakehouse}
объединяет достоинства обоих подходов: поверх дешёвого объектного хранилища строится
слой открытого табличного формата, обеспечивающий ACID-транзакции, эволюцию схемы,
моментальные снимки и параллельный доступ. К числу таких систем относятся Apache
Iceberg~\cite{iceberg-spec}, Delta Lake~\cite{armbrust2020delta}, Apache
Hudi~\cite{hudi-docs} и Apache Paimon~\cite{paimon-docs}.

При проектировании lakehouse-хранилища одним из ключевых решений является выбор
файлового формата, в котором физически хранятся данные. От него зависят степень
сжатия (а значит, объём хранилища и стоимость), скорость записи, скорость
аналитических запросов и эффективность пропуска нерелевантных данных. Де-факто
стандартом колоночного хранения в экосистеме больших данных является Apache Parquet;
широко применяется также Apache ORC. Однако оба формата разрабатывались более
десяти лет назад, их форматы кодирования отражают возможности процессоров того
времени и не используют современные приёмы --- векторизованное декодирование с
явной SIMD-параллельностью, адаптивное сжатие чисел с плавающей точкой, кодирование
строк по схожести префиксов. Эти приёмы реализованы в новом колоночном формате
Vortex~\cite{vortex-repo, vortex-docs}, разрабатываемом с прицелом на нативное
(написанное на Rust) векторизованное выполнение и переносимость метаданных. Вопрос
о том, в каких сценариях преимущества
Vortex проявляются на практике и можно ли интегрировать его в существующую
lakehouse-систему, на момент начала работы оставался открытым.

Актуальность работы определяется, во-первых, практической потребностью в
обоснованном выборе формата хранения при построении витрин данных на Apache Paimon,
а во-вторых, тем, что Vortex --- активно развивающийся формат, поддержка которого в
Apache Paimon отсутствовала, и интеграция вместе с экспериментальной оценкой
представляет самостоятельный интерес как для сообщества Apache Paimon, так и для
организации, в инфраструктуре которой выполнялась работа.

Цель работы --- определить область применимости колоночного формата Vortex
в lakehouse-системе Apache Paimon, разработать модуль его интеграции и выполнить
сравнительный анализ форматов хранения (Parquet, ORC, Vortex) на представительных
нагрузках.

Для достижения цели поставлены следующие задачи:

\begin{enumerate}
    \item провести анализ подходов к хранению аналитических данных --- от хранилищ
        данных (Hive Metastore, Impala) к архитектуре lakehouse и открытым табличным
        форматам;
    \item провести обзор колоночных файловых форматов (Parquet, ORC, Vortex, Lance)
        и обосновать исключение строчного формата Avro из сравнения;
    \item изучить архитектуру Apache Paimon --- LSM-дерево, таблицы с первичным
        ключом, слияние при чтении, --- и алгоритмы кодирования, сжатия и пропуска
        данных, применяемые форматами;
    \item разработать модуль интеграции файлового формата Vortex в Apache Paimon;
    \item провести функциональное тестирование полученного артефакта локально в
        контейнерах;
    \item развернуть на кластере окружение нагрузочного тестирования средствами
        Ansible;
    \item провести серию экспериментов и собрать метрики производительности;
    \item сформулировать рекомендации по применению форматов (опционально).
\end{enumerate}

Объект исследования --- колоночные файловые форматы хранения данных в
составе lakehouse-системы Apache Paimon. Предмет исследования --- алгоритмы
кодирования, сжатия и пропуска данных, реализуемые этими форматами, и их влияние на
производительность чтения и записи таблиц.

Методы исследования: анализ исходного кода и документации Apache Paimon,
Apache Flink и формата Vortex; разработка программного модуля на языках Java и Rust;
функциональное тестирование в контейнеризованном окружении; нагрузочное тестирование
на кластере с использованием стандартизированного набора запросов TPC-DS и
специально разработанного потокового сценария; статистическая обработка результатов
измерений (усреднение по повторным запускам, оценка разброса).

Практическая значимость. Разработанный модуль интеграции расширяет перечень
поддерживаемых Apache Paimon форматов хранения. Полученные количественные результаты
и рекомендации применимы при проектировании витрин данных и настройке
lakehouse-хранилищ. Выявленное архитектурное ограничение интерфейса форматов Apache
Paimon (отсутствие проталкивания агрегатов) указывает направление дальнейшего
развития системы.

Отчёт состоит из введения, трёх разделов, заключения, списка использованных
источников и приложений. Первый раздел содержит аналитический обзор: эволюцию систем
хранения, открытые табличные форматы, колоночные файловые форматы, архитектуру
Apache Paimon, алгоритмы кодирования и пропуска данных, а также постановку задачи.
Второй раздел посвящён разработке модуля интеграции формата Vortex, функциональному
тестированию, развёртыванию окружения нагрузочного тестирования, методике
экспериментов и обнаруженным в ходе работы проблемам. Третий раздел содержит
результаты экспериментов, их анализ по типам запросов и рекомендации по применению
форматов.
    % ВВЕДЕНИЕ

    % ── Раздел 1. Аналитический обзор ──
    \section{АНАЛИТИЧЕСКИЙ ОБЗОР}

В разделе рассматриваются эволюция систем хранения аналитических данных, открытые
табличные форматы как ключевой компонент архитектуры lakehouse, колоночные файловые
форматы хранения, архитектура системы Apache Paimon и применяемые форматами
алгоритмы кодирования, сжатия и пропуска данных. По итогам обзора формулируется
постановка задачи.

\subsection{Эволюция систем хранения аналитических данных}

Исторически аналитическая обработка данных опиралась на хранилища данных~\cite{inmon2005datawarehouse, kimball2013toolkit} (data
warehouse) --- специализированные реляционные системы, в которые данные загружаются
по расписанию через процессы ETL, приводятся к единой схеме «звезда» или
«снежинка» и затем обслуживают запросы бизнес-аналитики. Хранилище данных
обеспечивает строгую типизацию, управление схемой, транзакционность и
предсказуемую производительность запросов. Платой за это являются высокая стоимость
хранения (вычислительные узлы и хранилище связаны и масштабируются совместно),
жёсткость схемы (изменение требует перепроектирования ETL) и неспособность работать
со слабоструктурированными данными --- журналами, событиями, телеметрией.

С появлением распределённой файловой системы Hadoop (HDFS) сформировалась
альтернативная архитектура --- озеро данных (data lake). Данные складываются в HDFS
или объектное хранилище «как есть», в исходных форматах, а схема применяется во
время чтения (schema-on-read). Слой совместимости с SQL обеспечивался Hive Metastore~\cite{thusoo2009hive}
--- сервисом каталога, хранящим описания таблиц (расположение файлов, схему,
разбиение на разделы) в реляционной базе, --- и движками SQL поверх него: Apache
Hive~\cite{thusoo2009hive}, затем Apache Impala~\cite{kornacker2015impala}, Presto, Spark SQL. Озеро данных дёшево и масштабируемо,
позволяет хранить любые данные, но не даёт транзакционных гарантий: одновременная
запись и чтение, частичные обновления, атомарная замена набора файлов не
поддерживаются на уровне файловой системы. Hive-таблица, по сути, --- это директория
с файлами; согласованность при изменениях обеспечивается только дисциплиной на
стороне приложений.

Ограничения обоих подходов привели к появлению архитектуры lakehouse. Идея состоит
в том, чтобы сохранить дешёвое масштабируемое объектное хранилище озера данных, но
добавить поверх него слой метаданных, который превращает набор файлов в полноценную
таблицу с ACID-транзакциями, эволюцией схемы, моментальными снимками (snapshot
isolation), управлением версиями и параллельным доступом. Этот слой и называется
\textit{открытым табличным форматом}. Вычислительные движки (Spark, Flink, Trino и
другие) работают с таблицей через спецификацию табличного формата, а не напрямую с
файлами; разделение хранения и вычислений сохраняется, поскольку и данные, и
метаданные лежат в объектном хранилище. Современные реализации --- Apache Iceberg~\cite{iceberg-spec},
Delta Lake~\cite{armbrust2020delta}, Apache Hudi~\cite{hudi-docs} и Apache Paimon~\cite{paimon-docs} --- различаются деталями устройства
метаданных и набором поддерживаемых сценариев, но разделяют общую модель.

Особое место среди них занимает Apache Paimon. В отличие от Iceberg и Delta Lake,
ориентированных в первую очередь на таблицы только для добавления (append-only) и
пакетную загрузку, Paimon изначально проектировался под потоковую обработку и
таблицы с первичным ключом: в его основе лежит LSM-дерево, что делает эффективными
сценарии захвата изменений (Change Data Capture) и построения near-real-time-витрин.
Именно поэтому Apache Paimon выбран целевой системой настоящей работы.

    \subsection{Открытые табличные форматы}

Открытый табличный формат определяет, как из набора файлов данных, лежащих в
объектном хранилище, собирается логическая таблица. Минимальный набор обязанностей
табличного формата включает: ведение списка файлов, составляющих текущее состояние
таблицы; атомарную фиксацию изменений (добавление и удаление файлов одной
транзакцией); изоляцию снимков (читатель видит согласованное состояние на момент
начала чтения, не затрагиваемое параллельной записью); хранение схемы и её
эволюцию; статистики по файлам для пропуска нерелевантных данных; разбиение на
разделы. Рассмотрим основные реализации.

\textbf{Apache Iceberg}~\cite{iceberg-spec} ведёт метаданные в виде дерева: файл метаданных таблицы
ссылается на список манифестов (manifest list), каждый манифест перечисляет файлы
данных с их статистиками. Фиксация транзакции --- это атомарная замена указателя на
новый файл метаданных. Iceberg поддерживает скрытое разбиение на разделы (partition
evolution), эволюцию схемы без перезаписи данных, ветвление и теги снимков. Удаления
реализуются через позиционные и равенственные delete-файлы. Iceberg --- наиболее
широко принятый в индустрии формат, поддержанный всеми крупными движками.

\textbf{Delta Lake}~\cite{armbrust2020delta} использует журнал транзакций --- упорядоченную последовательность
JSON-файлов, каждый из которых описывает добавленные и удалённые файлы данных;
периодически журнал уплотняется в контрольную точку в формате Parquet. Delta
поддерживает обновления и удаления через перезапись затронутых файлов
(copy-on-write) либо через deletion vectors. Тесно интегрирован с Apache Spark.

\textbf{Apache Hudi}~\cite{hudi-docs} ориентирован на инкрементальную обработку и upsert по ключу
записи. Поддерживает два типа таблиц: copy-on-write (обновления порождают перезапись
файлов) и merge-on-read (обновления пишутся в журнальные файлы, состояние
вычисляется при чтении). Ведёт временную шкалу (timeline) операций и индекс,
сопоставляющий ключи записей файлам.

\textbf{Apache Paimon}~\cite{paimon-docs} (изначально Flink Table Store) выделяется тем, что состояние
таблицы с первичным ключом организовано как LSM-дерево, размещённое в объектном
хранилище. Свежие изменения сбрасываются в файлы уровня~0, фоновое слияние
перемещает их на нижележащие уровни; при чтении актуальная версия строки получается
объединением версий из всех уровней по полю последовательности. Метаданные ведутся
аналогично Iceberg: снимок ссылается на манифесты, манифесты --- на файлы данных со
статистиками. Такое устройство делает Paimon эффективным для потоковой записи,
захвата изменений и near-real-time-витрин --- сценариев, в которых таблица
постоянно обновляется небольшими порциями, а не перезаписывается целиком. Apache
Paimon поддерживает несколько файловых форматов хранения данных, выбираемых
параметром таблицы: Apache Parquet (по умолчанию), Apache ORC и Apache Avro;
добавление поддержки формата Vortex составляет практическую часть настоящей работы.

Для задач настоящей работы важно, что табличный формат и файловый формат --- разные
уровни абстракции: первый управляет составом и версиями набора файлов, второй ---
двоичным представлением данных внутри одного файла. Производительность чтения и
записи определяется обоими уровнями; в работе фиксируется табличный формат (Apache
Paimon) и варьируется файловый формат.

    \input{contents/2-03-file-formats}
    \subsection{Архитектура Apache Paimon}

Apache Paimon~\cite{paimon-docs} --- открытый табличный формат, объединяющий слой метаданных в стиле
Apache Iceberg с организацией данных в виде LSM-дерева~\cite{oneil1996lsm, luo2020lsm} в объектном хранилище. Ниже
рассматриваются модель данных, устройство таблицы с первичным ключом, путь чтения и
точка подключения файлового формата.

\textbf{Модель данных и метаданные.} Таблица Paimon --- это директория в объектном
хранилище, содержащая поддиректории данных и метаданных. Каждая фиксация (commit)
порождает новый \textit{снимок} (snapshot) --- небольшой файл, ссылающийся на список
\textit{манифестов}. Манифест перечисляет файлы данных, добавленные или удалённые
данной фиксацией, вместе со статистиками по столбцам каждого файла (минимум,
максимум, число null). Чтение всегда выполняется относительно конкретного снимка,
что обеспечивает изоляцию: параллельная запись создаёт новые снимки, не затрагивая
тот, который читает потребитель. Старые снимки и недостижимые файлы удаляются по
политике хранения (retention).

\textbf{Разбиение и бакеты.} Таблица может быть разбита на разделы (partitions) по
значениям выбранных столбцов; внутри раздела данные распределяются по фиксированному
числу \textit{бакетов} (buckets) хешированием ключа. Бакет --- единица параллелизма
записи и единица LSM-дерева: каждый бакет имеет собственное LSM-дерево. Запись в
бакет однопоточна, что устраняет конкуренцию за одно дерево; число бакетов задаёт
максимальный параллелизм записи.

\textbf{Таблица с первичным ключом и LSM-дерево.} В таблице с первичным ключом строки
идентифицируются ключом; обновления и удаления не перезаписывают данные на месте, а
добавляются как новые версии. Внутри бакета новые записи буферизуются в памяти в
отсортированном по ключу виде и при фиксации (по чекпойнту вычислительного движка)
сбрасываются на диск одним отсортированным неизменяемым файлом --- это \textit{уровень~0}
(\(L_0\)) LSM-дерева. Файлов \(L_0\) накапливается по одному на чекпойнт. Фоновая
операция \textit{слияния} (compaction) объединяет несколько отсортированных серий
(sorted runs) в файлы нижележащих уровней (\(L_1\), \(L_2\), \dots), отбрасывая
перезаписанные и помеченные на удаление записи и укрупняя файлы; слияние не
блокирует запись. Параметры \texttt{num-sorted-run.compaction-trigger} (число серий,
при котором запускается слияние) и \texttt{num-sorted-run.stop-trigger} (порог
обратного давления на запись) управляют его агрессивностью; \texttt{target-file-size}
задаёт целевой размер укрупнённых файлов. Схематически устройство показано на
рисунке~\ref{fig:lsm}.

\begin{figure}
\centering
\begin{tikzpicture}[
    font=\small,
    box/.style={draw, rounded corners=2pt, minimum height=7mm, minimum width=14mm, align=center},
    lvl/.style={font=\small\bfseries},
    >={Stealth[length=2mm]}
  ]
  % memory buffer
  \node[box, fill=black!5] (mem) at (0,3.4) {буфер в памяти\\(сортировка по ключу)};
  % L0
  \node[lvl] at (-4.6,2.2) {$L_0$};
  \node[box, fill=black!8] (l0a) at (-3.2,2.2) {файл};
  \node[box, fill=black!8] (l0b) at (-1.6,2.2) {файл};
  \node[box, fill=black!8] (l0c) at (0.0,2.2) {файл};
  \node at (1.4,2.2) {$\dots$};
  % L1
  \node[lvl] at (-4.6,1.0) {$L_1$};
  \node[box, fill=black!12, minimum width=22mm] (l1a) at (-2.8,1.0) {файл};
  \node[box, fill=black!12, minimum width=22mm] (l1b) at (-0.2,1.0) {файл};
  % L2
  \node[lvl] at (-4.6,-0.2) {$L_2$};
  \node[box, fill=black!16, minimum width=34mm] (l2a) at (-2.0,-0.2) {файл};
  \node[box, fill=black!16, minimum width=34mm] (l2b) at (2.0,-0.2) {файл};
  % arrows
  \draw[->] (mem) -- node[right]{\footnotesize flush по чекпойнту} (l0c);
  \draw[->] (l0a) -- ($(l1a.north)+(-0.4,0)$);
  \draw[->] (l0b) -- ($(l1a.north)+(0.4,0)$);
  \draw[->, dashed] (l0c) -- (l1b.north);
  \node[font=\footnotesize] at (3.3,1.55) {слияние ($L_0\!\to\!L_1$)};
  \draw[->] (l1a) -- (l2a.north);
  \draw[->] (l1b) -- ($(l2a.north)+(0.8,0)$);
  \node[font=\footnotesize] at (3.6,0.45) {слияние ($L_1\!\to\!L_2$)};
  % read side
  \node[box, fill=black!5, minimum width=30mm] (rd) at (5.6,3.0) {итератор слияния\\при чтении};
  \draw[->, thick] (l0c.east) to[out=0,in=180] ($(rd.west)+(0,0.3)$);
  \draw[->, thick] (l1b.east) to[out=0,in=180] (rd.west);
  \draw[->, thick] (l2b.east) to[out=0,in=180] ($(rd.west)+(0,-0.3)$);
\end{tikzpicture}
\caption{Устройство LSM-дерева бакета в Apache Paimon: буферизация в памяти, сброс в~\(L_0\) по чекпойнту, фоновое слияние на нижележащие уровни, объединение версий итератором слияния при чтении}
\label{fig:lsm}
\end{figure}

\textbf{Поле последовательности и вид изменения.} Чтобы итератор слияния выбирал
актуальную версию строки, таблица настраивается на поле последовательности
(\texttt{sequence.field}) --- монотонно возрастающую величину (например, момент
изменения), по максимуму которой определяется последняя версия ключа; поле вида
изменения (\texttt{rowkind.field}) различает вставку, обновление и удаление. Эта
схема естественно соответствует потокам захвата изменений (CDC).

\textbf{Путь чтения.} Чтение бакета порождает итератор слияния (merge iterator),
который параллельно читает отсортированные серии всех уровней и выдаёт строки в
порядке ключа, оставляя для каждого ключа версию с максимальным значением поля
последовательности. Если файлов \(L_0\) накопилось много (слияние не успело за
записью), итератор вынужден открыть и слить больше файлов --- стоимость чтения растёт.
Перед открытием файла применяется пропуск данных по статистикам из манифеста: если
диапазоны значений столбцов файла не пересекаются с предикатом запроса, файл
пропускается целиком; внутри файла дополнительно работает пропуск блоков средствами
самого файлового формата. Декодированные данные передаются вычислительному движку
пакетами строк (vectorized batches).

\textbf{Точка подключения файлового формата.} Файловый формат подключается к Paimon
через интерфейс \texttt{FileFormat} (модуль \texttt{paimon-common}), который
предоставляет фабрику записи (\texttt{FormatWriterFactory}) и фабрику чтения
(\texttt{FormatReaderFactory}). Фабрика записи создаёт по пути в хранилище объект
\texttt{FormatWriter}, принимающий строки и сериализующий их в файл выбранного
формата. Фабрика чтения создаёт по пути, проекции столбцов и (опционально)
предикату объект \texttt{FormatReaderFactory.FormatReader}, выдающий итератор пакетов
строк. Реализация формата также объявляет, какие предикаты она способна принять для
проталкивания, и формирует объект статистик по записанному файлу для манифеста.
Существующие реализации --- \texttt{paimon-format} (Parquet, ORC, Avro); добавление
модуля \texttt{paimon-vortex}, реализующего этот интерфейс поверх нативной библиотеки
Vortex, и составляет практическую часть работы.

\textbf{Интеграция с вычислительным движком.} В работе используется Apache Flink~\cite{carbone2015flink}:
потоковая запись --- через приёмник Paimon, чтение --- через источник Paimon в
пакетном режиме. Запись управляется чекпойнтами Flink: каждый чекпойнт фиксирует
снимок Paimon. Интервал чекпойнтов задаёт частоту фиксаций: чем он больше, тем
крупнее (и реже) файлы \(L_0\), но тем больше данных буферизуется и тем заметнее
всплеск ввода-вывода при сбросе. Этот параметр существенен для near-real-time-сценария
и отдельно варьируется в экспериментах.

    \subsection{Кодирование, сжатие и пропуск данных}

Производительность колоночного формата определяется тремя группами механизмов:
кодированием значений (компактное двоичное представление, использующее структуру
данных), пропуском нерелевантных данных (отказ от чтения и декодирования блоков,
заведомо не нужных запросу) и векторизованным декодированием (распаковка плотных
массивов значений с использованием параллелизма процессора). Рассмотрим их.

Базовые кодеки. Словарное кодирование (dictionary encoding): для
столбца с небольшим числом различных значений строится словарь, а сами значения
заменяются индексами в нём; выгодно для категориальных строк (страна, тип
устройства, метод запроса). Кодирование длин серий (run-length encoding,
RLE): последовательность одинаковых подряд идущих значений заменяется парой
«значение, длина серии»; выгодно для отсортированных или малоизменчивых столбцов.
Упаковка в биты (bit-packing): целые числа из ограниченного диапазона
хранятся в минимально необходимом числе бит, а не в полном машинном слове.
Кодирование относительно опорного значения (frame of reference, FOR): из
всех значений вычитается минимум блока, остатки бит-упаковываются~\cite{zukowski2006superscalar, lemire2015decoding}. Дельта-кодирование:
хранятся разности соседних значений; выгодно для монотонных последовательностей
(идентификаторы, временные метки). Поверх кодеков применяется блочное сжатие общего
назначения (Snappy, ZSTD, LZ4). Parquet и ORC используют этот арсенал; различие в
деталях (размеры блоков, выбор кодека, наличие фильтров Блума).

FastLanes. Vortex применяет для целочисленных данных схему бит-упаковки
FastLanes~\cite{afroozeh2023fastlanes}, спроектированную под SIMD-декодирование. Обычная бит-упаковка требует при
распаковке битовых сдвигов и масок с зависимостями по данным, что плохо ложится на
векторные инструкции. FastLanes переставляет биты значений в памяти («транспонирует»
их по полосам фиксированной ширины) так, чтобы распаковка \(k\) значений выполнялась
небольшим числом векторных операций без ветвлений и без зависимостей между полосами.
В сочетании с FOR и дельта-кодированием это даёт высокую пропускную способность
декодирования целочисленных столбцов.

ALP. Для чисел с плавающей точкой Vortex применяет кодек ALP~\cite{afroozeh2024alp} (Adaptive
Lossless floating-Point). Идея: многие вещественные значения в реальных данных ---
это десятичные дроби с небольшим числом знаков (координаты, цены, измерения),
которые при умножении на подходящую степень десяти становятся целыми. ALP подбирает
по выборке блока показатели степени, представляет основную массу значений как целые
числа (далее --- FOR и бит-упаковка), а немногочисленные значения, не укладывающиеся
в эту модель (исключения), хранит отдельно в исходном виде. Преобразование
обратимо без потерь. На столбцах с «настоящими» произвольными double ALP
вырождается, но на типичных прикладных данных существенно выигрывает у побайтового
хранения.

FSST. Для строк Vortex применяет кодек FSST~\cite{boncz2020fsst} (Fast Static Symbol Table). По
выборке строит таблицу из не более чем \num{255} коротких подстрок («символов»,
длиной до 8 байт), часто встречающихся в данных, и кодирует строки
последовательностью однобайтовых ссылок на символы (байт \texttt{0xFF} экранирует
буквальный байт). Декодирование --- простая подстановка по таблице. FSST особенно
эффективен на строках с регулярной структурой и общими префиксами: URL и пути
(\texttt{/api/v2/cat\_017}), идентификаторы (\texttt{sess\_0000001234},
\texttt{req\_0000000042}), доменные имена, --- то есть на данных, которые в
аналитических таблицах встречаются повсеместно. Существенно: эффективность FSST (и
прочих кодеков) зависит от порядка строк внутри блока. В отсортированном
блоке соседние строки похожи, символы переиспользуются плотно; в перемешанном блоке
выигрыш меньше. Это объясняет, почему в LSM-дереве выгоды Vortex полнее проявляются
на уплотнённых (отсортированных по ключу) данных, чем на свежих файлах \(L_0\).

Каскадирование кодеков. В Vortex кодеки --- композируемые: к уже
закодированному массиву применяется следующий кодек (например, словарь, затем
бит-упаковка индексов словаря, затем FOR). Выбор каскада адаптивен --- по выборке
оценивается, какая комбинация даёт наилучшее сжатие при приемлемой стоимости
декодирования.

Пропуск данных по статистикам. Файловый формат хранит для каждого блока
(группа строк / страница в Parquet, полоса / группа в ORC, блок в Vortex) статистики
--- минимум, максимум, число null. Если предикат запроса (\texttt{col BETWEEN a AND
b}, \texttt{col = v}, \texttt{col IS NOT NULL}) несовместим с диапазоном блока, блок
не читается. Та же информация на уровне всего файла дублируется в манифесте Paimon,
что позволяет отсечь файл ещё до его открытия. Тонкость: статистики по строковым
столбцам часто хранятся усечёнными (truncate, например, до 16 байт первого
несовпадающего символа) ради компактности; для столбцов, у которых различающая
часть выходит за границу усечения, пропуск перестаёт работать. В работе используется
режим полных статистик (\texttt{metadata.stats-mode = full}), чтобы исключить этот
эффект из сравнения.

Проталкивание предикатов. Помимо пропуска целых блоков, формат может
проталкивать предикат в декодер: вычислять условие на закодированных или
частично декодированных данных и выдавать только подходящие строки, не материализуя
остальные. Vortex поддерживает проталкивание предикатов на уровне нативной
библиотеки: Paimon-предикат преобразуется конвертером в нативное выражение Vortex
(операции сравнения \(=, \ne, <, \le, >, \ge\), проверка на null, конъюнкция и
дизъюнкция для всех поддерживаемых типов), и фильтрация выполняется в Rust-коде до
пересечения границы JNI. Это даёт выигрыш на запросах с селективным предикатом:
через границу JVM/native передаётся лишь малая доля строк.

Отсутствие проталкивания агрегатов. Принципиальное ограничение: интерфейс
файловых форматов Apache Paimon (\texttt{FileFormat}) не предусматривает
проталкивание агрегатов --- операции \texttt{SUM}, \texttt{COUNT}, \texttt{COUNT
DISTINCT}, \texttt{GROUP BY} выполняются вычислительным движком (Flink) над уже
декодированными данными в JVM. Это архитектурное ограничение самого Paimon, а не
конкретного формата. Для Parquet, у которого весь путь --- от декодера до движка
агрегации --- находится в JVM, граница JNI отсутствует. Для Vortex данные сначала
пересекают границу JNI (хотя и без копирования, через Arrow C Data Interface), а
затем агрегируются в JVM; на запросах вида «полное сканирование + агрегация без
селективного предиката» это даёт Vortex дополнительные накладные расходы. Устранение
ограничения --- проталкивание агрегатов в нативный SIMD-слой Vortex --- лежит за
рамками работы и обозначено как направление развития Apache Paimon.

Векторизованное декодирование. И Parquet, и ORC, и Vortex выдают данные
движку не построчно, а пакетами (vectorized batches) --- плотными массивами значений
столбцов, к которым движок применяет векторизованные операции. Vortex дополнительно
строит декодированные пакеты в раскладке Apache Arrow~\cite{arrow-cdata}, что устраняет промежуточное
копирование при передаче в Arrow-совместимые компоненты. Однако в Paimon результат
всё равно адаптируется к внутреннему представлению строк Flink, поэтому полностью
исключить накладные расходы границы не удаётся.

    \subsection{Постановка задачи}

Проведённый обзор позволяет уточнить задачу работы. Apache Paimon --- lakehouse-система
с LSM-деревом, поддерживающая несколько файловых форматов; среди современных
колоночных форматов выделяется Vortex с нативным векторизованным декодированием,
специализированными кодеками (FastLanes, ALP, FSST) и проталкиванием предикатов, но
его поддержка в Apache Paimon отсутствует. Требуется определить, в каких сценариях
Vortex даёт выигрыш, реализовать его интеграцию и измерить характеристики.

Конкретизация задач.

\begin{enumerate}
    \item Разработка модуля интеграции. Реализовать модуль
        \texttt{paimon-vortex}: фабрики записи и чтения поверх нативной библиотеки
        Vortex (доступ через JNI), конвертер Paimon-предикатов в нативные выражения
        Vortex, формирование статистик записанных файлов для манифеста, поддержку
        проекции столбцов и многопоточного чтения.
    \item Функциональное тестирование. Проверить корректность артефакта
        локально в контейнерах: запись и чтение таблиц с первичным ключом и таблиц
        только для добавления, слияние версий, эволюция числа бакетов, согласованность
        с другими форматами на одинаковых данных.
    \item Развёртывание окружения нагрузочного тестирования. Средствами
        Ansible развернуть на кластере (Apache Flink, Apache Paimon, объектное
        хранилище) воспроизводимое окружение для запуска нагрузочных задач, сбора
        метрик и управления жизненным циклом кластера.
    \item Эксперименты. Провести две серии измерений: (а)~batch-аналитика
        на наборе TPC-DS масштаба \SI{100}{\giga\byte} (таблицы только для добавления,
        крупные файлы) --- сравнение Parquet и Vortex по времени выполнения запросов;
        (б)~near-real-time-сценарий --- таблица с первичным ключом, потоковая запись
        с захватом изменений, параллельные пакетные аналитические запросы во время
        записи и после неё --- сравнение Parquet, ORC и Vortex по пропускной
        способности записи и времени запросов разных типов. Результаты статистически
        обработать.
    \item Рекомендации. По результатам сформулировать рекомендации по выбору
        файлового формата хранения для различных классов нагрузок в Apache Paimon.
\end{enumerate}

Гипотезы, проверяемые в работе:

\begin{enumerate}
    \item на batch-аналитике с селективными предикатами и крупными файлами Vortex
        быстрее Parquet за счёт нативного проталкивания предикатов и
        специализированных кодеков;
    \item на потоковой near-real-time-нагрузке Vortex не уступает Parquet по
        пропускной способности записи;
    \item на запросах с агрегацией без селективного предиката Vortex проигрывает
        Parquet из-за накладных расходов границы JVM/native в отсутствие
        проталкивания агрегатов;
    \item поведение Vortex под конкурентным чтением во время потоковой записи
        отличается от Parquet/ORC ввиду различия в реализации доступа к объектному
        хранилищу.
\end{enumerate}

Критерии оценки --- время выполнения запросов (среднее по повторным запускам и
выборкам, с учётом разброса), пропускная способность записи (строк в секунду), число
успешно завершённых запросов под нагрузкой; для качественных выводов ---
устойчивость метрик (наличие или отсутствие катастрофических выбросов).


    % ── Раздел 2. Разработка и развёртывание ──
    \section{РАЗРАБОТКА МОДУЛЯ ИНТЕГРАЦИИ И РАЗВЁРТЫВАНИЕ ОКРУЖЕНИЯ}

В разделе описаны разработанный модуль интеграции формата Vortex в Apache Paimon,
функциональное тестирование артефакта в контейнерах, развёртывание окружения
нагрузочного тестирования на кластере средствами Ansible, методика проведённых
экспериментов и проблемы, выявленные и решённые в ходе работы.

\subsection{Модуль интеграции формата Vortex в Apache Paimon}

Модуль реализован как набор подмодулей Maven в составе Apache Paimon (интеграция
выполнена в ответвлении репозитория, на основе версии Apache Paimon 1.4):
\texttt{paimon-vortex-jni} --- обёртки нативных методов и загрузчик библиотеки;
\texttt{paimon-vortex-format} --- реализация интерфейса \texttt{FileFormat} и
связанных классов; нативная библиотека Vortex собирается из исходных текстов на
Rust и упаковывается в артефакт. Логически модуль состоит из четырёх частей: путь
записи, путь чтения, конвертер предикатов и слой доступа к нативной библиотеке.

\textbf{Слой доступа к нативной библиотеке.} Класс \texttt{NativeFileMethods}
объявляет \texttt{native}-методы: открытие файла по URI с параметрами хранилища
(\texttt{open}), получение числа строк и схемы (\texttt{rowCount}, \texttt{dtype}),
запуск сканирования с проекцией столбцов, сериализованным предикатом и диапазоном
строк (\texttt{scan}), закрытие (\texttt{close}), а также операции уровня каталога
(перечисление и удаление файлов). Загрузчик \texttt{NativeLoader} извлекает
библиотеку под текущую платформу из ресурсов артефакта и подгружает её. Передача
больших массивов данных через границу JNI выполняется без копирования: нативная
сторона экспортирует декодированный пакет в раскладке Apache Arrow через Arrow C
Data Interface, а JVM-сторона импортирует его в \texttt{VectorSchemaRoot}.

\textbf{Путь записи.} \texttt{VortexFileFormat} создаёт \texttt{FormatWriterFactory},
которая по пути в хранилище порождает писатель: строки Paimon (\texttt{InternalRow})
преобразуются в пакеты Arrow заданного размера (\texttt{write.batch-size}) и
передаются нативному писателю Vortex; по завершении файла формируется объект
статистик по столбцам (минимум, максимум, число null) для записи в манифест Paimon.
Размер пакета записи существенен: он же определяет размер пакета, возвращаемого при
чтении одним нативным вызовом, поэтому от него зависит число пересечений границы JNI
при последующем сканировании. В работе используется увеличенный относительно
умолчания размер пакета, что снижает накладные расходы JNI при чтении (см.
подраздел~\ref{subsec:issues}).

\textbf{Путь чтения.} \texttt{VortexFileFormat} создаёт \texttt{FormatReaderFactory};
её метод порождает \texttt{VortexRecordsReader} по пути файла, проекции столбцов,
(опционально) предикату и параметрам хранилища. Читатель открывает нативный файл,
запускает сканирование (передавая список проецируемых столбцов, сериализованный
предикат и --- при наличии --- список номеров строк для выборочного чтения) и
получает итератор нативных массивов. Метод \texttt{readBatch} извлекает очередной
нативный массив, экспортирует его в Arrow \texttt{VectorSchemaRoot} (с
переиспользованием ранее выделенной структуры), преобразует средствами
\texttt{ArrowBatchReader} в итератор \texttt{InternalRow} и возвращает движку
вместе с информацией о позиции строки в файле (необходима Paimon для
deletion-векторов и слияния). Нативные ресурсы (массив, аллокатор Arrow, итератор,
дескриптор файла) освобождаются при закрытии читателя или при освобождении пакета.

\textbf{Конвертер предикатов.} \texttt{VortexPredicateConverter} преобразует
предикат Paimon в нативное выражение Vortex. Поддержаны операции сравнения
(\(=, \ne, <, \le, >, \ge\)), проверка на null, конъюнкция и дизъюнкция для всех
типов данных, для которых это имеет смысл (логический, целочисленные, строковый,
десятичный, дата, временная метка). Непереводимые предикаты не проталкиваются ---
фильтрация по ним остаётся на стороне движка; корректность при этом не нарушается.
Выражение сериализуется и передаётся нативной библиотеке, которая применяет его при
сканировании до пересечения границы JNI, что сокращает объём данных, передаваемых в
JVM, на запросах с селективным предикатом.

\textbf{Многопоточность.} Параллелизм обеспечивается на уровне Paimon и Flink: чтение
таблицы разбивается на разделённые задачи (splits) --- по бакетам и файлам, ---
каждая обрабатывается своим читателем в отдельном потоке (subtask) вычислительного
движка; запись в бакет однопоточна, параллелизм записи равен числу бакетов. Внутри
одного читателя нативное сканирование Vortex также может использовать несколько
потоков. Объекты модуля, разделяемые между задачами (фабрики), не имеют изменяемого
состояния; объекты, привязанные к одной задаче (читатель, аллокатор Arrow,
нативные дескрипторы), потоконебезопасны и используются строго из одного потока, что
соответствует контракту интерфейсов Paimon. Параметры памяти настраиваются с учётом
того, что нативная сторона использует прямые (off-heap) буферы вне кучи JVM ---
этому посвящён отдельный пункт подраздела~\ref{subsec:issues}.

Последовательность взаимодействия компонентов при чтении таблицы Paimon с форматом
Vortex показана на рисунке~\ref{fig:seq-read}.

\begin{figure}
\centering
\resizebox{\textwidth}{!}{%
\begin{tikzpicture}[
    font=\footnotesize,
    >={Stealth[length=2mm]},
    head/.style={draw, rounded corners=1.5pt, align=center, minimum height=8mm, minimum width=20mm, fill=black!4},
    msg/.style={->, thick},
    ret/.style={->, dashed},
  ]
  % --- координаты времени жизни ---
  \def\xA{0}      % Flink
  \def\xB{3.6}    % Paimon FileFormat
  \def\xC{7.2}    % VortexFileFormat
  \def\xD{10.8}   % VortexRecordsReader
  \def\xE{14.4}   % Vortex (Rust) via JNI
  \def\ytop{0}
  \def\ybot{-10.6}
  % --- головы участников ---
  \node[head] (A) at (\xA,\ytop) {Flink\\(движок)};
  \node[head] (B) at (\xB,\ytop) {Apache Paimon\\\texttt{FileFormat}};
  \node[head] (C) at (\xC,\ytop) {\texttt{VortexFile}\\\texttt{Format}};
  \node[head] (D) at (\xD,\ytop) {\texttt{VortexRecords}\\\texttt{Reader}};
  \node[head] (E) at (\xE,\ytop) {Vortex (Rust)\\через JNI};
  % --- линии жизни ---
  \foreach \x in {\xA,\xB,\xC,\xD,\xE} {
    \draw[dashed, black!55] (\x,-0.45) -- (\x,\ybot);
  }
  % --- сообщения ---
  % 1
  \draw[msg] (\xA,-1.2) -- node[above]{\texttt{createReaderFactory(проекция, предикат)}} (\xB,-1.2);
  % 2
  \draw[msg] (\xB,-1.9) -- node[above]{\texttt{createReaderFactory(\dots)}} (\xC,-1.9);
  % 3 self: convert predicate
  \draw[msg] (\xC,-2.6) -- ++(1.0,0) -- ++(0,-0.5) -- node[right, xshift=2pt]{\texttt{convert()} — предикат $\to$ нативное выражение} ++(-1.0,0);
  % 4 return factory
  \draw[ret] (\xC,-3.7) -- node[above]{\texttt{FormatReaderFactory}} (\xB,-3.7);
  % 5 new reader
  \draw[msg] (\xB,-4.4) -- node[above]{\texttt{new VortexRecordsReader(путь, тип, предикат)}} (\xD,-4.4);
  % 6 open + scan
  \draw[msg] (\xD,-5.1) -- node[above]{\texttt{open(uri)}; \texttt{scan(столбцы, предикат)}} (\xE,-5.1);
  \draw[ret] (\xE,-5.7) -- node[above]{итератор нативных массивов} (\xD,-5.7);
  % --- loop frame ---
  \draw[black!60] (-1.5,-6.2) rectangle (15.9,-9.3);
  \node[draw, fill=white, font=\scriptsize\bfseries, anchor=north west] at (-1.5,-6.2) {loop [пока есть пакеты]};
  % loop messages
  \draw[msg] (\xA,-6.9) -- node[above]{\texttt{readBatch()}} (\xD,-6.9);
  \draw[msg] (\xD,-7.6) -- node[above]{\texttt{next()}; \texttt{exportToArrow()}} (\xE,-7.6);
  \draw[ret] (\xE,-8.2) -- node[above]{Arrow \texttt{VectorSchemaRoot} (без копирования)} (\xD,-8.2);
  \draw[ret] (\xD,-8.9) -- node[above]{пакет строк (\texttt{InternalRow})} (\xA,-8.9);
  % --- close ---
  \draw[msg] (\xA,-9.8) -- node[above]{\texttt{close()}} (\xD,-9.8);
  \draw[msg] (\xD,-10.4) -- node[above]{освобождение нативных ресурсов} (\xE,-10.4);
\end{tikzpicture}%
}
\caption{Диаграмма последовательности: чтение таблицы Apache Paimon с файловым форматом Vortex}
\label{fig:seq-read}
\end{figure}

    \subsection{Функциональное тестирование в контейнерах}

Перед нагрузочными экспериментами корректность модуля проверена локально в
контейнеризованном окружении. Окружение разворачивается через Docker Compose и
включает: контейнер с Apache Flink (менеджер заданий и менеджер задач), контейнер с
объектным хранилищем, совместимым с протоколом S3 (MinIO), и контейнер-клиент для
запуска заданий. Артефакт Paimon с модулем \texttt{paimon-vortex} собирается из
исходных текстов; нативная библиотека Vortex кросс-компилируется под целевую
платформу контейнера.

Проверялись следующие свойства.

\begin{enumerate}
    \item \textbf{Запись и чтение таблицы только для добавления.} Создание таблицы с
        форматом Vortex, загрузка данных, чтение, сверка числа строк и контрольных
        агрегатов с эталоном.
    \item \textbf{Таблица с первичным ключом и слияние версий.} Начальная загрузка,
        затем серия обновлений по случайным ключам с возрастающим полем
        последовательности; чтение должно вернуть для каждого ключа последнюю версию.
        Корректность сверялась с результатом той же последовательности операций на
        формате Parquet.
    \item \textbf{Проекция столбцов и проталкивание предикатов.} Запросы с выборкой
        части столбцов и с селективными предикатами разных типов (целочисленные
        диапазоны, равенство строк, проверка на null); сверка результатов с эталоном
        и проверка по журналам, что предикат действительно протолкнут.
    \item \textbf{Эволюция числа бакетов и слияние.} Изменение числа бакетов с
        перераспределением данных; принудительное слияние; чтение после слияния.
    \item \textbf{Граничные случаи кодеков.} Столбцы с высокой и низкой
        кардинальностью, с большой долей null, со строками-префиксами, с числами с
        плавающей точкой --- проверка, что кодеки Vortex (FastLanes, ALP, FSST)
        активируются и данные восстанавливаются без потерь.
\end{enumerate}

В ходе тестирования выявлен и устранён ряд проблем сборки и конфигурации (несовместимость
версий Java и Flink на клиентском процессе, идемпотентность повторных запусков,
требования к памяти вне кучи JVM для нативных буферов Vortex, особенности
кросс-компиляции нативной библиотеки под версию системной библиотеки C целевой
платформы); они подробно рассмотрены в подразделе~\ref{subsec:issues}. Также
установлено существенное ограничение: целевой файловой системой для нативной
библиотеки Vortex является объектное хранилище, совместимое с S3; распределённая
файловая система Hadoop (HDFS) и схема URI \texttt{hdfs://} не поддерживаются, ---
что определило выбор объектного хранилища в качестве бэкенда и в контейнерном
окружении, и на кластере. После устранения проблем все перечисленные проверки
пройдены: артефакт признан функционально корректным и пригодным для нагрузочного
тестирования.

    \subsection{Развёртывание окружения нагрузочного тестирования}
\label{subsec:deployment}

Нагрузочные эксперименты проводились на кластере из пяти узлов. Программный стек:
Apache Flink (версия 2.x) как вычислительный движок, Apache Paimon (версия 1.4) с
модулем \texttt{paimon-vortex} как табличный формат, Java~17 в качестве среды
исполнения, объектное хранилище, совместимое с S3 (на базе Apache Ozone~\cite{ozone-docs}), как бэкенд
хранения. На каждом узле запускаются два менеджера задач (TaskManager) по 30 слотов
--- итого 10 менеджеров задач и 300 слотов в кластере; на одном из узлов
дополнительно работает менеджер заданий (JobManager). Решение запускать два менеджера
задач на узел вместо одного с 60 слотами принято осознанно: меньший размер кучи на
процесс (\SI{100}{\giga\byte} вместо \SI{200}{\giga\byte}) сокращает паузы полной
сборки мусора, а изоляция процессов означает, что пауза одного менеджера задач не
блокирует второй; это особенно важно для формата Parquet, чей путь декодирования и
агрегации полностью находится в куче JVM, при высоком параллелизме.

Развёртывание, обновление конфигурации и управление жизненным циклом кластера
автоматизированы средствами Ansible~\cite{ansible-docs}: набор ролей и плейбуков (\texttt{start.yml},
\texttt{stop.yml}, \texttt{config-update.yml}) разворачивает Flink на узлах,
раскладывает зависимости (в том числе артефакт Paimon с нативной библиотекой
Vortex), формирует конфигурацию из шаблонов и запускает или останавливает менеджеры
заданий и задач. Отдельный плейбук (\texttt{submit-job.yml}) собирает задание,
копирует его и зависимости на узел менеджера заданий, запускает через клиент Flink,
извлекает идентификатор задания и опрашивает его состояние до завершения, после чего
забирает с узла файл метрик. Поскольку доступ к хранилищу на стенде защищён
Kerberos, добавлен плейбук обновления билетов (\texttt{krb-renew.yml}), выполняемый
первым во всех остальных плейбуках. Запуск нагрузочной серии для нескольких форматов
оформлен сценарием-обёрткой, который последовательно прогоняет задание для каждого
формата с предварительной очисткой состояния и принудительной остановкой
параллельных заданий.

Схема развёртывания и взаимодействия компонентов при проведении эксперимента
приведена на рисунке~\ref{fig:deployment}.

\begin{figure}
\centering
\resizebox{\textwidth}{!}{%
\begin{tikzpicture}[
    font=\small,
    >={Stealth[length=2.4mm]},
    node distance=6mm,
    comp/.style={draw, rounded corners=2pt, align=center, minimum height=9mm, fill=black!4, inner sep=4pt},
    tm/.style={draw, rounded corners=1.5pt, align=center, minimum height=6mm, minimum width=14mm, fill=black!8, font=\footnotesize},
    nodebox/.style={draw, dashed, rounded corners=3pt, inner sep=6pt},
    flow/.style={->, thick},
    dataflow/.style={->, thick, dash pattern=on 3pt off 2pt},
  ]
  % --- Dev machine + Ansible ---
  \node[comp] (dev) {Рабочая станция\\(Ansible: \texttt{start}, \texttt{stop},\\\texttt{config-update}, \texttt{submit-job})};
  % --- JobManager node ---
  \node[comp, right=22mm of dev, yshift=14mm] (jm) {Узел~1: JobManager\\(драйвер задания,\\файл метрик CSV)};
  % --- TM nodes ---
  \node[tm, right=30mm of jm, yshift=22mm] (tm1a) {TaskManager};
  \node[tm, below=2.5mm of tm1a] (tm1b) {TaskManager};
  \node[nodebox, fit=(tm1a)(tm1b), label={[font=\footnotesize]above:Узел~1}] (n1) {};

  \node[tm, below=10mm of tm1b] (tm2a) {TaskManager};
  \node[tm, below=2.5mm of tm2a] (tm2b) {TaskManager};
  \node[nodebox, fit=(tm2a)(tm2b), label={[font=\footnotesize]above:Узел~2}] (n2) {};

  \node[font=\large] (dots) at ($(tm2b)+(0,-1.0)$) {$\vdots$};

  \node[tm, below=14mm of tm2b] (tm5a) {TaskManager};
  \node[tm, below=2.5mm of tm5a] (tm5b) {TaskManager};
  \node[nodebox, fit=(tm5a)(tm5b), label={[font=\footnotesize]below:Узел~5}] (n5) {};

  \node[comp, below=24mm of jm, xshift=2mm] (s3) {Объектное хранилище\\(S3 / Apache Ozone):\\данные Paimon, метаданные,\\чекпойнты Flink};

  % --- arrows ---
  \draw[flow] (dev) -- node[above, sloped, font=\footnotesize]{ssh / Ansible} (jm);
  \draw[flow] (dev.south) to[out=-60, in=180] node[below, font=\footnotesize, pos=0.35]{развёртывание, конфигурация} (s3.west);
  \draw[flow] (jm.east) -- node[above, sloped, font=\footnotesize]{планирование задач} (n1.west);
  \draw[flow] (jm.east) to[out=0, in=180] (n2.west);
  \draw[flow] (jm.east) to[out=-25, in=180] (n5.west);
  \draw[dataflow] (n1.south) to[out=-90, in=20] node[right, font=\footnotesize, pos=0.4]{чтение / запись} (s3.east);
  \draw[dataflow] (n2.west) to[out=180, in=70] (s3.north east);
  \draw[dataflow] (n5.south west) to[out=200, in=-30] (s3.south east);
  \draw[flow] (s3.west) to[out=160, in=-100] node[left, font=\footnotesize, pos=0.6]{файл метрик} (jm.south);
  \draw[flow] (jm.north) to[out=120, in=0] node[above, font=\footnotesize, pos=0.6]{выгрузка результатов} (dev.east);
\end{tikzpicture}%
}
\caption{Схема развёртывания окружения нагрузочного тестирования: автоматизация через Ansible, кластер из пяти узлов (по два менеджера задач на узел), объектное хранилище как бэкенд данных и чекпойнтов}
\label{fig:deployment}
\end{figure}

    \subsection{Методика экспериментов}
\label{subsec:methodology}

Проведены две серии экспериментов, соответствующие двум основным сценариям
использования Apache Paimon: пакетная аналитика над таблицами только для добавления и
near-real-time-витрина на таблице с первичным ключом при потоковой записи.

\subsubsection{Сценарий 1: пакетная аналитика на наборе TPC-DS}

TPC-DS~\cite{tpcds-spec} --- стандартизированный набор для оценки систем поддержки принятия решений:
схема розничной торговли (таблицы фактов \texttt{store\_sales}, \texttt{web\_sales},
\texttt{catalog\_sales} и связанные измерения), генератор данных с управляемым
масштабом и \num{99} шаблонов запросов (с учётом вариантов --- \num{103} запроса),
покрывающих сканирование, фильтрацию, соединения, группировку и оконные функции.
Использован масштаб \SI{100}{\giga\byte}; данные сгенерированы, загружены в Apache
Paimon как таблицы только для добавления, разбиты на разделы (типичный размер файла
раздела --- порядка \SI{256}{\mega\byte}), параллелизм запросов --- 300. Сравнивались
форматы Parquet и Vortex по времени выполнения каждого запроса; для каждого запроса
выполнено несколько повторов, бралось характерное (устойчивое) значение. Запись в
этом сценарии не измерялась: загрузка крупных наборов упиралась в пропускную
способность объектного хранилища, что не позволяло корректно сопоставить форматы по
записи (этот сценарий покрывается второй серией).

\subsubsection{Сценарий 2: near-real-time-витрина при потоковой записи}

Разработано задание Flink, моделирующее near-real-time-витрину. Используется таблица
с первичным ключом, полем последовательности и полем вида изменения (схема захвата
изменений). Схема таблицы --- журнал событий из 25 столбцов, подобранный так, чтобы
дать форматам возможность проявить разнообразие кодеков: категориальные строки
низкой кардинальности (страна, тип устройства, операционная система, метод запроса
--- словарное и RLE-кодирование), строки с регулярной структурой и общими
префиксами (идентификаторы сессии и запроса, путь URL, версия ОС, строка
user-agent --- FSST), монотонные и узкодиапазонные целые (код ответа, задержка,
ширина экрана --- FOR и бит-упаковка), числа с плавающей точкой (координаты, выручка
--- ALP), столбцы с большой долей null (идентификатор родительского запроса ---
\SI{30}{\percent}, идентификатор объявления --- \SI{50}{\percent}, выручка ---
\SI{90}{\percent}). Ширина в 25 столбцов сопоставима с таблицами фактов TPC-DS ---
реалистичная форма аналитической таблицы. Состав схемы приведён в
таблице~\ref{tab:schema}.

\begin{table}
    \caption{Схема таблицы near-real-time-витрины (25 столбцов)}
    \label{tab:schema}
    {\small
    \begin{tabularx}{\textwidth}{|>{\raggedright\arraybackslash}p{0.22\textwidth}|>{\raggedright\arraybackslash}p{0.40\textwidth}|X|}\hline
    \textbf{Группа} & \textbf{Столбцы} & \textbf{Ожидаемые кодеки} \\ \hline
    Ключ и служебные & \texttt{id} (первичный ключ), \texttt{seq} (последовательность), \texttt{op} (вид изменения) & дельта / FOR; константа в фазе обновления \\ \hline
    Категории & \texttt{country\_code}, \texttt{device\_type}, \texttt{os\_name}, \texttt{event\_type}, \texttt{request\_method}, \texttt{referrer\_host} & словарь, RLE \\ \hline
    Префиксные строки & \texttt{session\_id}, \texttt{request\_id}, \texttt{parent\_req\_id}, \texttt{url\_path}, \texttt{user\_agent}, \texttt{os\_version} & FSST \\ \hline
    Целые & \texttt{user\_id}, \texttt{status\_code}, \texttt{latency\_ms}, \texttt{bytes\_sent}, \texttt{screen\_width}, \texttt{ad\_id} & FOR, бит-упаковка, дельта \\ \hline
    С плавающей точкой & \texttt{geo\_lat}, \texttt{geo\_lon}, \texttt{revenue} & ALP \\ \hline
    Временная метка & \texttt{ts} & дельта \\ \hline
    \end{tabularx}
    }
\end{table}

Эксперимент состоит из четырёх фаз. INIT --- начальная загрузка \num{50000000} строк
с последовательными идентификаторами (поле последовательности равно нулю); измеряется
пропускная способность записи. UPDATE --- поток обновлений по случайным
идентификаторам из загруженного диапазона (поле последовательности возрастает) с
ограничением скорости (для получения воспроизводимой нагрузки скорость зафиксирована
на уровне порядка 100\,000~строк/с; объём подобран так, чтобы фаза длилась
сопоставимое время в разных запусках). Интервал чекпойнтов Flink установлен в
240~с --- это режим near-real-time: фиксации происходят раз в несколько минут,
файлы уровня~0 крупные. DURING --- во время фазы обновления параллельно, с
интервалом в несколько минут, запускаются пакетные аналитические запросы по текущему
снимку (несколько выборок); это моделирует потребителя, читающего витрину во время её
наполнения. FINAL --- после завершения записи те же запросы выполняются по
итоговому снимку. Различие фаз DURING и FINAL принципиально: в фазе DURING читатель
конкурирует с активной записью за ввод-вывод, а LSM-дерево содержит много ещё не
слитых файлов уровня~0; в фазе FINAL запись остановлена, а дерево относительно
упорядочено (хотя полное принудительное слияние перед итоговыми запросами не
выполняется --- читается состояние «как есть после остановки записи», что и
соответствует реальной витрине).

Набор аналитических запросов охватывает характерные типы нагрузки; он приведён в
таблице~\ref{tab:queries}. Каждый запрос материализует результат в таблицу-приёмник
(чтобы исключить влияние клиента на измерение); метрикой служит время выполнения в
миллисекундах. Весь цикл (создание таблицы --- INIT --- UPDATE с DURING-выборками ---
FINAL) повторяется несколько раз (число повторов --- не менее трёх); состояние между
повторами сбрасывается. По повторам и DURING-выборкам вычисляются как среднее, так
и медиана; основной метрикой принята медиана, так как DURING-сэмплы могут
попадать в момент массового сброса данных по чекпойнту Flink (в эти моменты writer
одновременно пишет в объектное хранилище и данные таблицы, и состояние чекпойнта,
размещаемое в той же директории хранилища), что приводит к единичным «хвостовым»
сэмплам длительностью в десятки секунд; такие сэмплы инфлируют среднее, но не
влияют на медиану. Среднее приводится отдельно как индикатор подверженности
формата выбросам.

\begin{xltabular}{\textwidth}{|>{\raggedright\arraybackslash}p{0.40\textwidth}|X|}
    \caption{Аналитические запросы near-real-time-сценария}\label{tab:queries} \\ \hline
    \textbf{Запрос} & \textbf{Содержание и проверяемый механизм} \\ \hline
    \endfirsthead
    \caption*{Продолжение таблицы~\ref{tab:queries}} \\ \hline
    \textbf{Запрос} & \textbf{Содержание и проверяемый механизм} \\ \hline
    \endhead
    \hline
    \endfoot
    \hline
    \endlastfoot
    \texttt{POINT} & выборка одной строки по первичному ключу; точечный доступ, фильтры Блума и статистики \\ \hline
    \texttt{FULLSCAN\_PROJ} & полное сканирование с проекцией; нет предиката --- проверяет «чистую» скорость декодирования и накладные расходы границы \\ \hline
    \texttt{RANGE\_SESSION} & диапазонный фильтр по \texttt{session\_id} (строка-префикс); FSST + пропуск файлов по статистикам строкового столбца \\ \hline
    \texttt{RANGE\_URL} & диапазонный фильтр по \texttt{url\_path}; то же на пути URL \\ \hline
    \texttt{AGG\_SUM\_BYTES} & \texttt{SUM} по числовому столбцу с полным сканированием; агрегация без селективного предиката \\ \hline
    \texttt{AGG\_COUNT\_DISTINCT} & \texttt{COUNT(DISTINCT)} по строковому столбцу; тяжёлая агрегация с полным сканированием \\ \hline
    \texttt{AGG\_FILTERED} & \texttt{COUNT} с диапазонным предикатом; предикат отсекает данные, агрегация амортизируется \\ \hline
    \texttt{AGG\_GROUP\_COUNTRY} & \texttt{GROUP BY} по категориальному столбцу; группировка по столбцу низкой кардинальности \\ \hline
    \texttt{AGG\_MULTIDIM} & многомерная группировка с несколькими агрегатами \\ \hline
    \texttt{FUNNEL} & многошаговый запрос (последовательность событий пользователя); несколько проходов и предикатов \\ \hline
    \texttt{JOIN\_USER\_ACTIVITY} & соединение по \texttt{user\_id} с агрегацией по типу события; column pruning при соединении \\ \hline
    \texttt{JOIN\_SESSION\_LATENCY} & соединение по \texttt{session\_id} с диапазонным фильтром по задержке \\ \hline
    \texttt{JOIN\_REVENUE\_PROFILE} & подзапрос-сборка активных пользователей и соединение с их строками покупок; 5 из 25 столбцов \\ \hline
\end{xltabular}

По итогам предварительных запусков зафиксирован
рабочий набор параметров: распределение слотов между записью и чтением (общая сумма
--- весь кластер); число бакетов; размер буфера записи под выбранный интервал
чекпойнтов; увеличенный размер пакета записи Vortex (снижает накладные расходы JNI
при чтении); полные статистики по столбцам (\texttt{metadata.stats-mode = full});
расширенные политики хранения снимков и манифестов (необходимы для устойчивости при
частых фиксациях и параллельном чтении); файловые форматы --- без специфичных тонких
настроек размера файлов и страниц. От ряда гипотез пришлось отказаться: использование
строчного формата Avro для свежих файлов уровня~0 ухудшает аналитику Vortex на
итоговом снимке; агрессивные настройки размера файлов и страниц благоприятствуют
Parquet и раздувают хранилище за счёт накопления удаляемых файлов в окне хранения;
эти результаты обсуждаются в подразделе~\ref{subsec:issues}.

Измерение пропускной способности записи требует оговорки. В фазе обновления все
строки имеют одинаковое значение поля последовательности, поэтому Paimon
дедуплицирует записи с совпадающим ключом и последовательностью в буфере записи до
сброса на диск. При большом числе обновлений на один ключ значительная доля строк не
доходит до хранилища, и измеренная «пропускная способность» в фазе обновления
отражает скорость генерации и дедупликации в памяти, а не скорость записи на диск.
Поэтому в выводах по записи используется в первую очередь пропускная способность фазы
INIT (последовательные ключи, реальная запись), а показатели фазы UPDATE
интерпретируются с этой оговоркой.

Однократный запуск даёт разброс результатов до 25--40\,\% между
идентичными конфигурациями: на результат влияют разогрев JIT-компилятора первого из
прогоняемых форматов, стоимость первой загрузки нативной библиотеки, размещение
процессов по узлам, состояние пулов соединений с хранилищем, накопленное состояние
кластера. Для статистически значимого сравнения требуется несколько повторов с
чередованием порядка форматов; единичные выбросы
интерпретировались как артефакты совпадения момента запроса с пиком ввода-вывода от
записи или паузой сборки мусора.

    \subsection{Проблемы, выявленные и решённые в ходе работы}
\label{subsec:issues}

В ходе разработки модуля, развёртывания окружения и проведения экспериментов
выявлен ряд проблем; их систематизация полезна как практический результат работы.
Проблемы сгруппированы по этапам.

\subsubsection{Сборка и локальное тестирование}

\begin{enumerate}
    \item \textbf{Совместимость Java~17, Flink и DataStream API на клиентском
        процессе.} Запуск задания падал с ошибкой доступа к закрытому полю модуля
        \texttt{java.base} при инициализации транслятора приёмника DataStream.
        Причина --- параметры \texttt{--add-opens} были заданы только для менеджеров
        задач, но не для клиентского процесса. Решение --- задать
        \texttt{env.java.opts.all} (применяется ко всем JVM: менеджер заданий,
        менеджеры задач, клиент).
    \item \textbf{Идемпотентность повторных запусков.} Повторный запуск поверх
        существующих файлов наращивал LSM-дерево и искажал метрики. Решение ---
        принудительная очистка директории таблицы в хранилище перед каждым запуском.
    \item \textbf{Память вне кучи JVM для нативных буферов Vortex.} Чтение Vortex с
        конфигурацией памяти по умолчанию приводило к ошибке исчерпания прямой
        (off-heap) памяти при выделении буферов Arrow/JNI; Parquet и ORC на той же
        конфигурации проходили. Вывод --- Vortex принципиально требует увеличенного
        бюджета off-heap-памяти; конфигурацию менеджеров задач необходимо настраивать
        с учётом нативной стороны.
    \item \textbf{Кросс-компиляция нативной библиотеки.} Сборка нативной библиотеки
        Vortex требовала более новой версии системной библиотеки C, чем доступна на
        целевой версии операционной системы кластера; решение --- кросс-компиляция
        под более старую версию (конфигурация зафиксирована в сборочных контейнерах).
    \item \textbf{Целевая файловая система.} Установлено, что нативная библиотека
        Vortex рассчитана на объектное хранилище, совместимое с S3; распределённая
        файловая система Hadoop (HDFS) и схема URI \texttt{hdfs://} не поддерживаются.
        Это определило выбор бэкенда хранения.
\end{enumerate}

\subsubsection{Развёртывание кластера и автоматизация}

\begin{enumerate}
    \item \textbf{Гонка между завершением быстрого задания и его проверкой.} Плейбук
        проверки ожидал появления задания в состоянии «выполняется», но быстрые
        пакетные задания успевали завершиться до первого опроса. Решение --- извлекать
        идентификатор задания из вывода клиента и опрашивать его конечное состояние,
        принимая в том числе «завершено», и отдельно сигнализировать только об ошибке.
    \item \textbf{Отсутствие классов Hadoop при десериализации графа задания.}
        Перезапуск кластера через обновление конфигурации не экспортировал путь к
        библиотекам Hadoop, в отличие от обычного запуска; решение --- добавить
        экспорт переменных окружения в задачи перезапуска менеджеров.
    \item \textbf{Самоуничтожение shell-обёрток при остановке процессов.} Команды
        вида \texttt{pkill -f 'TaskManagerRunner'} матчили собственную командную
        строку shell-обёртки и посылали сигнал самим себе. Решение --- сузить шаблон
        до полного имени класса и применить приём с символьным классом в регулярном
        выражении (\texttt{[T]askManagerRunner}), который соответствует процессу, но
        не буквальному тексту в командной строке обёртки.
    \item \textbf{Kerberos для доступа к хранилищу.} Для личной учётной записи на
        стенде недоступен режим keytab; пришлось использовать режим кэша билетов с
        отдельным плейбуком их продления, выполняемым первым во всех остальных
        плейбуках.
    \item \textbf{Запись файла метрик на узле менеджера заданий.} Драйвер задания
        выполняется на узле менеджера заданий, где не существует локально созданной
        директории для метрик; решение --- создавать родительский каталог
        непосредственно из кода задания.
    \item \textbf{Автоназначение портов для двух менеджеров задач на узле.} Два
        менеджера задач на одном узле регистрируются только при незаданном порте RPC
        (автоназначение); явное задание порта приводит к конфликту.
\end{enumerate}

\subsubsection{Хранилище: каталог Apache Paimon поверх Apache Ozone}

При повторном создании каталога возникала ошибка «путь должен быть директорией»:
объектное хранилище хранит «пустые директории» как нулевой объект с завершающим
слешем, и при последующем запросе слой совместимости с S3 принимал этот объект-маркер
за обычный файл. Решение --- удалять именно объект-маркер, не затрагивая дочерние
префиксы, после чего обращение к пути корректно разрешается как директория по
дочерним объектам; в перспективе --- использовать собственную схему URI хранилища или
размещать склад в корне бакета.

\subsubsection{Настройка потокового сценария}

\begin{enumerate}
    \item \textbf{Залипание генератора данных на параллелизме~1.} Без указания
        ограничения скорости генератор не шардировался по подзадачам, и пропускная
        способность залипала на пределе одного потока. Решение --- при отсутствии
        ограничения подставлять заведомо большой лимит, что активирует шардирование.
    \item \textbf{Потеря файлов под читателем при частых фиксациях.} Политики
        хранения снимков и манифестов по умолчанию при коротком интервале чекпойнтов
        агрессивно удаляли старые снимки; стартующий писатель или параллельный
        читатель обращались к уже удалённому манифесту. Решение --- расширенные
        политики хранения и асинхронный режим очистки.
    \item \textbf{Исчерпание прямой памяти при высоком параллелизме чтения.}
        Конфигурация с одним крупным менеджером задач на узел упиралась в лимит
        прямой памяти при большом параллелизме чтения и параллельной записи; решение
        --- два менеджера задач на узел меньшего размера (см.
        подраздел~\ref{subsec:deployment}).
    \item \textbf{Накладные расходы JNI при чтении Vortex.} При размере пакета по
        умолчанию (\num{1024} строк) полное сканирование бакета порождало тысячи
        пересечений границы JNI на файл. Решение --- увеличить размер пакета записи
        (он же определяет размер возвращаемого при чтении пакета) до десятков тысяч
        строк, что сократило число пересечений границы в десятки раз и устранило
        отставание Vortex по «чистому» сканированию.
    \item \textbf{Влияние принудительного полного слияния перед итоговыми запросами.}
        Попытка выполнять полное слияние перед фазой FINAL приводила к тому, что
        запросы стартовали до фактического завершения слияния (системная процедура
        слияния запускается как отдельное задание движка, а ожидание её результата
        возвращалось преждевременно); кроме того, агрессивное слияние конкурировало с
        читателем за ввод-вывод. От принудительного слияния отказались: фаза FINAL
        читает состояние «как есть после остановки записи», что и соответствует
        реальной витрине.
\end{enumerate}

\subsubsection{Отклонённые гипотезы по тонкой настройке}

\begin{enumerate}
    \item \textbf{Строчный формат для свежих файлов уровня~0.} Гипотеза «дешёвая
        построчная запись свежих дельт, колоночный формат --- только на нижних
        уровнях» дала прирост пропускной способности записи, но замедлила
        аналитические запросы Vortex на итоговом снимке (итератор слияния вынужден
        совмещать построчные дельты с колоночными файлами, нативный конвейер не
        активируется); при достаточном параллелизме кластера в этом приёме нет нужды.
    \item \textbf{Агрессивные настройки размера файлов и страниц.} Гипотеза
        «крупнее файлы --- лучше амортизация накладных расходов» подтвердилась для
        Parquet (его векторизованный читатель оптимален на крупных страницах), но не
        для Vortex; преимущества Vortex на ряде запросов при этом исчезли, а объём
        хранилища вырос из-за накопления удаляемых файлов в окне хранения при более
        частом слиянии. Использованы значения по умолчанию.
    \item \textbf{Агрессивное снижение порога слияния.} Снижение порога запуска
        слияния учащает фоновые слияния, что увеличивает конкуренцию за ввод-вывод с
        читателем в фазе DURING (наблюдались единичные многосекундные провалы времени
        запросов) и раздувает окно хранения снимков. Использованы значения по умолчанию
        (порог слияния, целевой размер файла).
\end{enumerate}

Что касается записи на наборе TPC-DS --- провести её измерение не удалось: загрузка
крупных наборов в хранилище упиралась в его пропускную способность, что не позволяло
корректно сопоставить форматы; этот аспект покрывается потоковым сценарием.


    % ── Раздел 3. Результаты экспериментов ──
    \section{РЕЗУЛЬТАТЫ ЭКСПЕРИМЕНТОВ И АНАЛИЗ}

В разделе приведены результаты двух серий экспериментов --- пакетной аналитики на
наборе TPC-DS и near-real-time-сценария при потоковой записи, --- их анализ по типам
запросов и рекомендации по применению форматов хранения.

\subsection{Сценарий 1: пакетная аналитика на наборе TPC-DS}

На наборе TPC-DS~\cite{tpcds-spec} масштаба \SI{100}{\giga\byte} сравнивались форматы Parquet и Vortex
по времени выполнения \num{103} запросов (таблицы только для добавления, крупные
файлы разделов, параллелизм 300). Из \num{103} запросов Vortex быстрее Parquet более
чем на \SI{10}{\percent} на \num{82} запросах; среднее ускорение по запросу ---
порядка \SI{25}{\percent}; пиковые ускорения достигают \SI{80}{\percent}. Parquet
быстрее Vortex лишь на пяти запросах. Характерные запросы приведены в
таблице~\ref{tab:tpcds}.

\begin{table}
    \caption{Характерные результаты на наборе TPC-DS (масштаб \SI{100}{\giga\byte}): относительное изменение времени выполнения Vortex по сравнению с Parquet}
    \label{tab:tpcds}
    \begin{tabularx}{\textwidth}{|l|c|X|}\hline
    \textbf{Запрос} & \textbf{Vortex к Parquet} & \textbf{Особенности запроса (объяснение)} \\ \hline
    q51 & $\approx -80\,\%$ & оконные нарастающие суммы по продажам, мало столбцов, селективные условия --- column pruning и нативное проталкивание предикатов работают в полную силу \\ \hline
    q42 & $\approx -52\,\%$ & агрегация продаж по категории с фильтром по дате --- селективный предикат + узкая проекция \\ \hline
    q35 & $\approx -52\,\%$ & демографический срез с подзапросами существования --- многократное сканирование немногих столбцов \\ \hline
    q73 & $\approx -41\,\%$ & подсчёт покупателей по числу покупок с фильтрами --- предикат отсекает данные до агрегации \\ \hline
    q88 & $\approx -35\,\%$ & восемь подсчётов с разными временными окнами --- восьмикратное сканирование одного столбца, выигрыш от column pruning умножается \\ \hline
    \multicolumn{3}{|c|}{\textit{запросы, на которых быстрее Parquet}} \\ \hline
    q1 & медленнее & «покупатели, вернувшие товары»: коррелированный подзапрос, агрегация возвратов --- результат во многом определяется стороной соединения и агрегации в JVM \\ \hline
    q2, q28 & медленнее & многократные объединения подзапросов с агрегацией без селективного предиката \\ \hline
    q83, q99 & медленнее & сводки возвратов и заказов с группировкой; преобладает агрегация над декодированными данными \\ \hline
    \end{tabularx}
\end{table}

\textbf{Анализ.} Условия сценария максимально благоприятны для Vortex: таблицы
только для добавления (нет слияния версий), крупные файлы (накладные расходы границы
JNI амортизируются на больших пакетах), реальные данные с избыточностью (кодеки
FastLanes, ALP, FSST дают высокую степень сжатия), запросы с селективными
предикатами и узкими проекциями (нативное проталкивание предикатов сокращает объём
данных, пересекающих границу JVM/native). Пиковое ускорение в \SI{80}{\percent} на
q51 и линейный рост выигрыша с числом сканирований немногих столбцов (q88 --- восемь
сканирований одного столбца) прямо подтверждают первую гипотезу работы.

Запросы, на которых быстрее Parquet, объединяет общая черта --- преобладание
агрегации без селективного предиката над относительно широкими наборами столбцов:
основная работа выполняется не в декодере, а в движке агрегации, который для Parquet
полностью находится в JVM, а для Vortex отделён границей JNI. Это согласуется с
третьей гипотезой и с обсуждавшимся в подразделе~\ref{subsec:methodology}
архитектурным ограничением --- отсутствием проталкивания агрегатов в интерфейсе
форматов Apache Paimon.

Измерить пропускную способность записи на этом наборе не удалось: загрузка крупных
объёмов упиралась в пропускную способность объектного хранилища, и сопоставление
форматов по записи было бы некорректным. Этот аспект покрывается второй серией
экспериментов, где запись выполняется потоком с ограничением скорости.

    \subsection{Сценарий 2: near-real-time-витрина при потоковой записи}

В сценарии сравнивались форматы ORC, Parquet и Vortex на таблице с первичным ключом
(схема журнала событий из 25 столбцов, таблица~\ref{tab:schema}) при потоковой записи
с захватом изменений и параллельных аналитических запросах. Ниже приведены результаты
представительного запуска (три повтора полного цикла, ограничение скорости фазы
обновления --- порядка 100\,000~строк/с, интервал чекпойнтов --- 240~с);
для пояснения поведения под конкурентным чтением привлекаются данные ещё одного
запуска.

\textbf{Запись.} Пропускная способность фазы начальной загрузки (последовательные
ключи, реальная запись на диск) практически одинакова у всех трёх форматов
(таблица~\ref{tab:write}); различие в пределах разброса. В фазе обновления показатели
также близки, но интерпретируются с оговоркой: при постоянном значении поля
последовательности значительная доля обновлений дедуплицируется в буфере записи и не
доходит до хранилища, поэтому величина отражает скорость генерации и дедупликации в
памяти. Вывод: Vortex не уступает Parquet и ORC по пропускной способности записи ---
вторая гипотеза работы подтверждается.

\begin{table}
    \caption{Пропускная способность записи, тыс. строк в секунду (среднее по повторам)}
    \label{tab:write}
    \begin{tabularx}{\textwidth}{|X|c|c|}\hline
    \textbf{Формат} & \textbf{Начальная загрузка (INIT)} & \textbf{Обновление (UPDATE)$^{*}$} \\ \hline
    ORC     & 1604 & 91 \\ \hline
    Parquet & 1532 & 98 \\ \hline
    Vortex  & 1527 & 91 \\ \hline
    \multicolumn{3}{|p{0.95\textwidth}|}{\footnotesize $^{*}$~значение фазы обновления отражает скорость генерации и дедупликации в буфере записи (постоянное поле последовательности), а не скорость записи на диск.} \\ \hline
    \end{tabularx}
\end{table}

\textbf{Чтение: усреднённые показатели.} Средние по всем запросам, выборкам и
повторам времена приведены в таблице~\ref{tab:read-avg}. В фазе FINAL (запись
остановлена) Vortex --- быстрейший в среднем; в фазе DURING (параллельная запись)
картина менее однозначна: ORC и Vortex близки, Parquet заметно медленнее за счёт
отдельных катастрофических выбросов (см. ниже).

\begin{table}
    \caption{Среднее время выполнения запроса, мс (по всем запросам, выборкам и повторам)}
    \label{tab:read-avg}
    \begin{tabularx}{\textwidth}{|X|c|c|}\hline
    \textbf{Формат} & \textbf{Фаза DURING} & \textbf{Фаза FINAL} \\ \hline
    ORC     & 7594  & 7633 \\ \hline
    Parquet & 9481  & 7767 \\ \hline
    Vortex  & 7712  & \textbf{6736} \\ \hline
    \end{tabularx}
\end{table}

\textbf{Чтение: фаза FINAL по запросам.} Подробные результаты приведены в
таблице~\ref{tab:read-final} (полужирным выделено наименьшее время по строке). Vortex
выигрывает на десяти запросах из тринадцати, проигрывая только многомерной агрегации
(\texttt{AGG\_MULTIDIM}, незначительно), диапазонному фильтру по URL
(\texttt{RANGE\_URL}) и одному из соединений (\texttt{JOIN\_REVENUE\_PROFILE}). Особенно
велик отрыв на запросах, чувствительных к упорядоченности данных и к качеству
кодирования строк-префиксов: \texttt{RANGE\_SESSION} (соединение с диапазонным
фильтром по идентификатору сессии) --- Vortex быстрее ORC примерно на \SI{20}{\percent};
\texttt{FUNNEL} (многошаговый запрос) --- примерно на \SI{24}{\percent};
\texttt{AGG\_FILTERED} (подсчёт с предикатом) --- примерно на \SI{24}{\percent}.

\begin{table}
    \caption{Время выполнения запросов в фазе FINAL, мс (среднее по повторам; полужирным --- наименьшее по строке)}
    \label{tab:read-final}
    \begin{tabularx}{\textwidth}{|X|c|c|c|}\hline
    \textbf{Запрос} & \textbf{ORC} & \textbf{Parquet} & \textbf{Vortex} \\ \hline
    \texttt{POINT}                & 3944 & 4735 & \textbf{3766} \\ \hline
    \texttt{FULLSCAN\_PROJ}        & 6633 & 5866 & \textbf{5819} \\ \hline
    \texttt{RANGE\_SESSION}        & 7191 & 7673 & \textbf{5801} \\ \hline
    \texttt{RANGE\_URL}            & 7392 & \textbf{5891} & 6451 \\ \hline
    \texttt{AGG\_SUM\_BYTES}       & 5807 & 5836 & \textbf{5781} \\ \hline
    \texttt{AGG\_COUNT\_DISTINCT}  & 8614 & 9184 & \textbf{7819} \\ \hline
    \texttt{AGG\_FILTERED}         & 9345 & 9226 & \textbf{7115} \\ \hline
    \texttt{AGG\_GROUP\_COUNTRY}   & 5943 & 6819 & \textbf{5804} \\ \hline
    \texttt{AGG\_MULTIDIM}         & 7410 & \textbf{6681} & 7288 \\ \hline
    \texttt{FUNNEL}                & 11807 & 11505 & \textbf{8929} \\ \hline
    \texttt{JOIN\_USER\_ACTIVITY}  & 8224 & 9518 & \textbf{8145} \\ \hline
    \texttt{JOIN\_SESSION\_LATENCY}& 9557 & 8926 & \textbf{7343} \\ \hline
    \texttt{JOIN\_REVENUE\_PROFILE}& \textbf{7369} & 9117 & 7512 \\ \hline
    \multicolumn{1}{|r|}{\textbf{выигрышей}} & 1 & 2 & \textbf{10} \\ \hline
    \end{tabularx}
\end{table}

\textbf{Чтение: фаза DURING по запросам.} Результаты приведены в
таблице~\ref{tab:read-during}. Здесь Vortex выигрывает на шести запросах, ORC --- на
пяти, Parquet --- на двух; однако ключевое наблюдение --- не число выигрышей, а
\textit{разброс}. ORC меняется между фазами DURING и FINAL крайне незначительно
(среднее 7594 против 7633\,мс); Parquet деградирует существенно (9481 против 7767\,мс)
и ловит катастрофические выбросы (например, \texttt{AGG\_COUNT\_DISTINCT} ---
24\,801~мс в фазе DURING против 9184~мс в FINAL). Vortex деградирует
умеренно. В другом запуске под более тяжёлой конкурентной нагрузкой выбросы наблюдались
и у Vortex: время полного сканирования и соединений в отдельных DURING-выборках
вырастало в разы --- это попадание запроса в момент крупного сброса данных по
чекпойнту (100\,000~строк/с $\times$ 240~с $=$ около 24~млн строк
одним сбросом --- значительный всплеск ввода-вывода в объектное хранилище).

\begin{table}
    \caption{Время выполнения запросов в фазе DURING, мс (среднее по выборкам и повторам; полужирным --- наименьшее по строке)}
    \label{tab:read-during}
    \begin{tabularx}{\textwidth}{|X|c|c|c|}\hline
    \textbf{Запрос} & \textbf{ORC} & \textbf{Parquet} & \textbf{Vortex} \\ \hline
    \texttt{POINT}                & 4009  & \textbf{3712}  & 3718 \\ \hline
    \texttt{FULLSCAN\_PROJ}        & 7127  & \textbf{5776}  & 6093 \\ \hline
    \texttt{RANGE\_SESSION}        & \textbf{7358}  & 13651 & 10035 \\ \hline
    \texttt{RANGE\_URL}            & 6647  & 6786  & \textbf{6161} \\ \hline
    \texttt{AGG\_SUM\_BYTES}       & 5820  & 5711  & \textbf{5695} \\ \hline
    \texttt{AGG\_COUNT\_DISTINCT}  & \textbf{8616}  & 24801 & 12805 \\ \hline
    \texttt{AGG\_FILTERED}         & 9157  & 9933  & \textbf{7371} \\ \hline
    \texttt{AGG\_GROUP\_COUNTRY}   & \textbf{6003}  & 7191  & 6518 \\ \hline
    \texttt{AGG\_MULTIDIM}         & \textbf{6577}  & 7264  & 6730 \\ \hline
    \texttt{FUNNEL}                & 12384 & 12135 & \textbf{10355} \\ \hline
    \texttt{JOIN\_USER\_ACTIVITY}  & \textbf{8045}  & 8636  & 8492 \\ \hline
    \texttt{JOIN\_SESSION\_LATENCY}& 8672  & 8902  & \textbf{8011} \\ \hline
    \texttt{JOIN\_REVENUE\_PROFILE}& 8302  & 8757  & \textbf{8277} \\ \hline
    \multicolumn{1}{|r|}{\textbf{выигрышей}} & 5 & 2 & 6 \\ \hline
    \end{tabularx}
\end{table}

\textbf{Промежуточные выводы по сценарию.} Vortex эквивалентен Parquet и ORC по
пропускной способности записи (гипотеза~2 подтверждена); на итоговом снимке Vortex ---
быстрейший формат в среднем и на большинстве запросов, особенно на запросах с
диапазонным фильтром по строке-префиксу и с предикатом перед агрегацией; на чистой
агрегации без селективного предиката Vortex иногда уступает (гипотеза~3 частично
подтверждена и здесь --- \texttt{AGG\_MULTIDIM}); под конкурентной записью все
форматы деградируют, причём ORC --- наименее (зрелый нативный читатель с буферизацией
ввода-вывода), Parquet --- наиболее, с катастрофическими выбросами в моменты сброса
данных по чекпойнту, Vortex --- умеренно (гипотеза~4 подтверждена в части различия
поведения форматов под конкурентным чтением).

    \subsection{Паттерны по типам запросов и архитектурные наблюдения}

Объединение результатов двух серий позволяет выделить устойчивые закономерности
(таблица~\ref{tab:patterns}). В потоковом сценарии сравнение ведётся по медиане как
менее зашумлённой метрике, отражающей типичную задержку, к которой и чувствителен
потребитель витрины данных.

\begin{table}
    \caption{Сравнение форматов по типам запросов и операций (NRT-сценарий, медиана)}
    \label{tab:patterns}
    \begin{tabularx}{\textwidth}{|>{\raggedright\arraybackslash}p{0.24\textwidth}|>{\raggedright\arraybackslash}p{0.20\textwidth}|>{\raggedright\arraybackslash}p{0.20\textwidth}|X|}\hline
    \textbf{Тип операции} & \textbf{Vortex против Parquet} & \textbf{Vortex против ORC} & \textbf{Объяснение} \\ \hline
    Запись (INIT) & сопоставимы & сопоставимы & оба буферизуют строки и сбрасывают по чекпойнту, различия в пределах разброса \\ \hline
    Точечный доступ по ключу (\texttt{POINT}) & сопоставимы (разница доли процента) & сопоставимы & фильтры Блума и статистики работают у всех; на маленьких выборках различия в пределах разброса \\ \hline
    Полное сканирование с проекцией (\texttt{FULLSCAN\_PROJ}) & сопоставимы (разница доли процента) & Vortex чуть быстрее в FINAL ($\approx 1\,\%$) & column pruning + плотные кодеки работают на упорядоченных данных; на свежих L0-файлах эффект слабее \\ \hline
    Диапазонный фильтр по строке-префиксу (\texttt{RANGE\_SESSION}, \texttt{RANGE\_URL}) & Vortex заметно быстрее (\texttt{RANGE\_SESSION} в DURING $\approx 3{\times}$ быстрее) & Vortex стабильно быстрее (\texttt{RANGE\_SESSION} в DURING $\approx 25\,\%$; \texttt{RANGE\_URL} $\approx 4\,\%$) & FSST-кодирование строк с общими префиксами + нативное проталкивание предиката + пропуск файлов по полным статистикам \\ \hline
    Подсчёт с селективным предикатом (\texttt{AGG\_FILTERED}) & Vortex быстрее ($\approx 21\,\%$ в FINAL) & Vortex стабильно быстрее ($\approx 21\,\%$ в FINAL, $\approx 4\,\%$ в DURING) & предикат отсекает данные до пересечения границы JNI; выигрыш от проталкивания преобладает над накладными расходами \\ \hline
    Многошаговая аналитика (\texttt{FUNNEL}, тяжёлые соединения с фильтром) & Vortex стабильно быстрее (\texttt{FUNNEL} $\approx 33\,\%$) & Vortex стабильно быстрее (\texttt{FUNNEL} $\approx 33\,\%$, \texttt{JOIN\_REVENUE\_PROFILE} $\approx 20\,\%$) & многократные проходы и предикаты, выигрыш от column pruning и проталкивания умножается \\ \hline
    Простая группировка/многомерная агрегация без предиката (\texttt{AGG\_GROUP\_COUNTRY}, \texttt{AGG\_MULTIDIM}) & Parquet иногда быстрее (единицы процентов) & ORC/Parquet иногда быстрее (единицы процентов) & чистая агрегация в JVM, для Vortex --- через границу JNI; нет проталкивания агрегатов в интерфейсе Paimon \\ \hline
    Хвосты задержки под конкурентной записью (среднее против медианы) & Vortex существенно у́же (среднее $\approx$ медиана, у Parquet среднее на $\approx 21\,\%$ выше медианы) & оба формата дают узкие хвосты (среднее $\approx$ медиана) разными способами & пул соединений Hadoop S3A эпизодически насыщается на сбросе чекпойнта; Vortex изолирован независимым нативным пулом, ORC --- агрессивной буферизацией C++-читателя \\ \hline
    \end{tabularx}
\end{table}

\textbf{Где Vortex реально лучше ORC.} Если выбрать устойчивые наблюдения и
оставить запросы, на которых разница превышает разброс, преимущество Vortex над ORC
сосредоточено в трёх классах запросов: запросы с селективным предикатом
(проталкивание в нативный декодер до границы JNI), диапазонные запросы по строкам с
регулярной структурой (FSST + пропуск по полным статистикам), многошаговая аналитика
с несколькими предикатами и проходами. На простой группировке по столбцу низкой
кардинальности и на многомерной агрегации без предиката ORC оказывается не хуже
или лучше, на точечных запросах --- сопоставим.

\textbf{Эффект упорядоченности данных.} По медиане ORC и Parquet практически не
меняют времени между фазами FINAL и DURING (отношение DURING/FINAL близко к единице),
а Vortex в DURING медленнее самого себя в FINAL на $\approx 22\,\%$. Это
прямое следствие зависимости специализированных кодеков Vortex (прежде всего FSST)
от \textit{порядка} строк внутри блока: в свежих файлах уровня~0, формируемых
по чекпойнту, данные ещё не отсортированы по ключу слиянием, и плотность
кодирования снижается; после остановки записи и работы фонового слияния файлы
сортированы по ключу, и кодеки работают в полную плотность. ORC и Parquet таким
эффектом не обладают --- их кодеки менее чувствительны к порядку, --- поэтому их
типичная задержка между фазами и не меняется.

\textbf{О выборе метрики усреднения и хвостах задержки.} В DURING-фазе отдельные
сэмплы чтения могут попадать в момент массового сброса данных по чекпойнту, когда
writer одновременно пишет в объектное хранилище и данные таблицы, и состояние
чекпойнта Flink (последнее размещается в той же директории объектного хранилища).
Среднее по сэмплам, особенно при небольшом числе повторов, чувствительно к таким
попаданиям. В работе в качестве основной метрики использована медиана. Соотнесение
среднего и медианы по DURING-фазе показывает: у ORC и Vortex среднее и медиана
практически совпадают (узкие хвосты задержки), у Parquet среднее на $\approx 21\,\%$
выше медианы --- среднее инфлировано несколькими крупными выбросами в десятки
секунд. Это согласуется с архитектурным предположением о пуле соединений: читатель
Parquet ходит в объектное хранилище через общий пул слоя совместимости Hadoop S3A
и в момент сброса данных по чекпойнту иногда простаивает в ожидании соединения, что
и даёт катастрофически медленную выборку. Читатель Vortex использует независимый
нативный пул и таких выбросов не даёт. Читатель ORC, формально использующий тот же
общий пул, в этом эксперименте также не даёт выбросов --- за счёт агрессивной
буферизации ввода-вывода в нативном C++-читателе, поглощающей эпизодические
задержки S3-операций.

\textbf{Почему Vortex доминирует на TPC-DS, но не доминирует на потоковом
LSM-сценарии.} На TPC-DS таблицы только для добавления, файлы крупные, путь чтения
прямой колоночный, накладные расходы границы JNI амортизируются на больших пакетах,
а запросы --- с высокой селективностью. На потоковом сценарии таблица --- с первичным
ключом и слиянием при чтении: итератор слияния открывает много файлов (особенно
свежих, не слитых файлов уровня~0), накладные расходы границы JNI накапливаются на
пер-файловой основе, данные в свежих файлах не упорядочены по ключу --- кодеки
Vortex (прежде всего FSST) теряют часть эффективности, --- а селективность запросов
в среднем ниже. Поэтому преимущество Vortex здесь скромнее и проявляется в первую
очередь на упорядоченном итоговом снимке и на запросах с селективным предикатом.
Для Vortex \textit{область наибольшей применимости} --- именно batch-аналитика на
таблицах только для добавления, не LSM с первичным ключом.

\textbf{Почему Parquet иногда быстрее на агрегации при полном сканировании.}
Проталкивание агрегатов отсутствует в интерфейсе файловых форматов Apache Paimon в
принципе --- это не специфика Vortex. Агрегаты вычисляются движком Flink над
декодированными данными. Для Parquet весь конвейер --- декодер и агрегатор --- в JVM,
границы нет. Для Vortex данные декодируются нативно, затем пересекают границу JNI
(без копирования, но с фиксированной стоимостью на пакет) и агрегируются в JVM; на
миллионах пакетов это даёт заметную добавку. Устранение ограничения --- проталкивание
агрегатов в нативный SIMD-слой Vortex --- естественное направление развития Apache
Paimon и потенциальный вклад в его экосистему.

\textbf{Почему два менеджера задач на узел лучше одного крупного.} Меньший размер
кучи на процесс ускоряет полную сборку мусора и сокращает паузы; изоляция процессов
означает, что пауза одного менеджера задач не блокирует второй; это особенно важно
для Parquet, чей путь декодирования и агрегации полностью в куче JVM, при высоком
параллелизме. Компактор Vortex использует память вне кучи и меньше зависит от размера
кучи, но также выигрывает от изоляции.

\subsection{Рекомендации по применению форматов хранения}

Рекомендации сформулированы с учётом сценария использования.

\textbf{1. Batch-аналитика, таблицы только для добавления.} Это \textit{основная
область применимости Vortex}. На наборе TPC-DS масштаба \SI{100}{\giga\byte} Vortex
быстрее Parquet на \num{82} запросах из \num{103}, среднее ускорение --- порядка
\SI{25}{\percent}, пиковое --- до \SI{80}{\percent} (запросы с селективным
предикатом, узкой проекцией, многократным сканированием немногих столбцов). Условия
этого сценария --- крупные файлы, отсутствие слияния версий, упорядоченные данные,
высокая селективность --- максимально благоприятны для специализированных кодеков и
нативного проталкивания предиката. \textbf{Рекомендуется выбирать Vortex} как
формат хранения данных для аналитических витрин на основе таблиц только для
добавления (фактовые таблицы, журналы, исторические данные, материализованные
выборки).

\textbf{2. Near-real-time-витрина (таблица с первичным ключом, потоковая запись).}
По типичной задержке (медиана) \textbf{Vortex} --- быстрейший формат в обеих фазах
и на большинстве запросов (девять из тринадцати); рекомендуется, если в нагрузке
преобладают запросы с селективным предикатом (отсечение до агрегации), диапазонные
запросы по строкам с регулярной структурой (идентификаторы, пути, URL) и
многошаговая аналитика (последовательности событий, соединения с диапазонными
фильтрами). \textbf{ORC} --- разумный выбор, если в нагрузке преобладают простая
группировка по столбцу низкой кардинальности, чистая многомерная агрегация без
предиката или если важна узкая задержка отдельных сэмплов (наихудший случай, p99) ---
ORC демонстрирует одинаково узкие хвосты с Vortex, а в отдельных запросах ему
способствует словарное кодирование. ORC также даёт самый компактный объём хранения
(\SI{17}{\giga\byte} против \SI{18}{\giga\byte} у Parquet и \SI{20}{\giga\byte} у
Vortex на потоковом сценарии; на TPC-DS преимущество ORC меньше). \textbf{Parquet}
--- разумный универсальный выбор по умолчанию: не уступает по записи, хуже по
медиане чтения, зрелый и широко поддержанный; основной недостаток --- эпизодические
выбросы времени запросов под конкурентной записью в моменты сброса данных по
чекпойнту (среднее заметно выше медианы).

\textbf{3. Чего не рекомендуется.} \textbf{Lance} в текущем состоянии интеграции с
Apache Paimon не рекомендуется для потоковых таблиц с первичным ключом (требуется
доработка слияния версий; на OLAP-сценарии отстаёт от остальных форматов в 2--3
раза); его область --- нагрузки машинного обучения с произвольным доступом к строкам
и векторным поиском, выходящие за рамки настоящей работы. \textbf{Apache Avro} ---
строчный формат, в принципе непригоден для OLAP-нагрузки; применение его для
свежих файлов уровня~0 в LSM-дереве было проверено и отклонено --- аналитические
запросы Vortex на итоговом снимке замедляются за счёт совмещения построчных дельт
с колоночными файлами в итераторе слияния.

\textbf{4. Конфигурация (общие соображения для NRT-витрин).} Распределять слоты
между записью и чтением так, чтобы читатель не блокировался; увеличить размер
пакета записи Vortex до десятков тысяч строк (снижает накладные расходы JNI при
чтении в десятки раз); включить полные статистики по столбцам
(\texttt{metadata.stats-mode = full}); настроить расширенные политики хранения
снимков и манифестов под частые фиксации; не применять агрессивных тонких
настроек размера файлов и страниц (благоприятствуют Parquet и раздувают
хранилище) и не использовать строчный формат для свежих файлов уровня~0;
увеличивать интервал чекпойнтов следует с учётом того, что это укрупняет всплеск
ввода-вывода при сбросе.


    \conclusion

В работе решена задача определения области применимости колоночного формата Vortex в
lakehouse-системе Apache Paimon, разработки модуля его интеграции и сравнительного
анализа форматов хранения. Поставленные задачи выполнены.

\begin{enumerate}
    \item Проведён аналитический обзор эволюции систем хранения аналитических данных
        --- от хранилищ данных и Hive Metastore к архитектуре lakehouse и открытым
        табличным форматам (Iceberg, Delta Lake, Hudi, Paimon); обоснован выбор
        Apache Paimon как целевой системы (LSM-дерево, ориентация на потоковую запись
        и таблицы с первичным ключом).
    \item Выполнен обзор колоночных файловых форматов (Parquet, ORC, Vortex, Lance);
        обосновано исключение строчного формата Avro из сравнения; разобраны
        алгоритмы кодирования, сжатия и пропуска данных --- словарное и RLE-кодирование,
        бит-упаковка, FOR и дельта-кодирование, FastLanes, ALP, FSST, статистики и
        фильтры Блума, проталкивание предикатов; выявлено архитектурное ограничение
        интерфейса форматов Apache Paimon --- отсутствие проталкивания агрегатов.
    \item Изучена архитектура Apache Paimon --- метаданные (снимки, манифесты),
        разбиение и бакеты, LSM-дерево таблицы с первичным ключом, слияние, путь
        чтения через итератор слияния, точка подключения файлового формата, интеграция
        с Apache Flink.
    \item Разработан и интегрирован в Apache Paimon модуль поддержки формата Vortex:
        реализация интерфейса \texttt{FileFormat} (фабрики записи и чтения),
        \texttt{VortexRecordsReader}, конвертер предикатов
        \texttt{VortexPredicateConverter}, слой доступа к нативной библиотеке через
        Java Native Interface с передачей данных без копирования через Arrow C Data
        Interface, формирование статистик записанных файлов, поддержка проекции
        столбцов и многопоточного чтения.
    \item Проведено функциональное тестирование артефакта локально в контейнерах:
        запись и чтение таблиц с первичным ключом и таблиц только для добавления,
        слияние версий, эволюция числа бакетов, проекция столбцов и проталкивание
        предикатов, граничные случаи кодеков; выявлен и устранён ряд проблем сборки и
        конфигурации.
    \item Развёрнуто на кластере из пяти узлов окружение нагрузочного тестирования
        средствами Ansible (Apache Flink, Apache Paimon с модулем
        \texttt{paimon-vortex}, объектное хранилище); автоматизированы развёртывание,
        обновление конфигурации, запуск заданий, сбор метрик, управление жизненным
        циклом кластера.
    \item Проведены две серии экспериментов --- batch-аналитика на наборе TPC-DS
        масштаба \SI{100}{\giga\byte} и near-real-time-сценарий с потоковой записью и
        параллельными аналитическими запросами; собраны и статистически обработаны
        метрики; результаты проанализированы по типам запросов.
    \item Сформулированы рекомендации по выбору файлового формата хранения для
        различных классов нагрузок в Apache Paimon.
\end{enumerate}

Гипотезы работы подтверждены: на batch-аналитике с селективными предикатами и
крупными файлами Vortex превосходит Parquet (среднее ускорение порядка
\SI{25}{\percent}, в пике --- до \SI{80}{\percent}); на потоковой
near-real-time-нагрузке Vortex эквивалентен Parquet и ORC по пропускной способности
записи и устойчивее под конкурентным чтением; на запросах с агрегацией без
селективного предиката Vortex иногда уступает Parquet --- следствие отсутствия
проталкивания агрегатов в интерфейсе форматов Apache Paimon.

Практическая значимость работы: расширен перечень поддерживаемых Apache Paimon
форматов хранения; получены количественные результаты и рекомендации, применимые при
проектировании витрин данных и настройке lakehouse-хранилищ; систематизированы
проблемы развёртывания и эксплуатации Apache Paimon с нативным файловым форматом на
кластере с объектным хранилищем.

Направления дальнейшей работы: проталкивание агрегатов в нативный SIMD-слой Vortex
(устранение выявленного архитектурного ограничения и потенциальный вклад в
экосистему Apache Paimon); поддержка дополнительных операций проталкивания и типов
данных в конвертере предикатов; исследование поведения форматов при ещё более высоком
параллелизме и иных профилях нагрузки; доработка интеграции формата Lance для
таблиц с первичным ключом.
      % ЗАКЛЮЧЕНИЕ

    \printbibliography                    % СПИСОК ИСПОЛЬЗОВАННЫХ ИСТОЧНИКОВ

    \appendix                             % ПРИЛОЖЕНИЯ
    \appendixsection{Исходные тексты проекта}

Все исходные тексты, разработанные и использованные в рамках работы, размещены в
публичной GitHub-организации
\href{https://github.com/orgs/mai-diploma-kochkozharov/repositories}{\texttt{mai-diploma-kochkozharov}}
(\url{https://github.com/orgs/mai-diploma-kochkozharov/repositories}).
Организация содержит следующие репозитории.

\begin{itemize}
    \item \textbf{\href{https://github.com/mai-diploma-kochkozharov/paimon}{\texttt{paimon}}}
        (ветка \texttt{vortex-on-1.4.0-rc1}). Ответвление репозитория
        \texttt{apache/paimon} на базе версии~1.4. Содержит разработанные в
        рамках работы подмодули: \texttt{paimon-vortex-jni} (обёртки нативных
        методов и загрузчик библиотеки), \texttt{paimon-vortex-format}
        (реализация интерфейса \texttt{FileFormat} --- фабрики записи и чтения,
        \texttt{VortexFormatWriter}, \texttt{VortexRecordsReader}, конвертер
        предикатов \texttt{VortexPredicateConverter}), а также интегрированную
        сборку нативной библиотеки Vortex.

    \item \textbf{\href{https://github.com/mai-diploma-kochkozharov/paimon-flink-job}{\texttt{paimon-flink-job}}}
        (ветка \texttt{main}). Нагрузочные задания Apache Flink~+ Apache Paimon:
        задание потокового near-real-time-сценария
        \texttt{StreamingWriteBenchmarkJob}, задание пакетной записи
        \texttt{WriteBenchmarkJob}, исполнитель стандартизованного набора TPC-DS
        \texttt{TpcdsBenchmarkJob}, сценарий-обёртка многоформатного прогона и
        Maven-конфигурация сборки.

    \item \textbf{\href{https://github.com/mai-diploma-kochkozharov/flink-lang-deploy}{\texttt{flink-lang-deploy}}}
        (ветка \texttt{master}). Развёртывание кластера Apache Flink средствами
        Ansible: роли установки и конфигурации, плейбуки
        \texttt{start.yml}, \texttt{stop.yml}, \texttt{config-update.yml},
        \texttt{submit-job.yml}, \texttt{krb-renew.yml}, шаблоны конфигурации
        Apache Flink, инвентарь кластера.

    \item \textbf{\href{https://github.com/mai-diploma-kochkozharov/paimon-vortex-deploy}{\texttt{paimon-vortex-deploy}}}
        (ветка \texttt{glibc-2.28}). Сборка артефактов в контейнерах с
        кросс-компиляцией нативной библиотеки Vortex под целевую версию
        системной библиотеки C кластера; конфигурация Docker Compose локального
        окружения для функциональных smoke-тестов; Makefile-цели сборки и
        запуска заданий.

    \item \textbf{\href{https://github.com/mai-diploma-kochkozharov/vortex-jni-test}{\texttt{vortex-jni-test}}}
        (ветка \texttt{master}). Изолированный проект для тестирования слоя
        доступа к нативной библиотеке Vortex через Java Native Interface:
        проверка корректности экспорта данных в Apache Arrow без копирования,
        граничные случаи кодеков, отладка нативных вызовов вне основного
        конвейера Paimon.

    \item \textbf{\href{https://github.com/mai-diploma-kochkozharov/paimon-benchmark-results}{\texttt{paimon-benchmark-results}}}
        (ветка \texttt{master}). Результаты экспериментов и сценарии их обработки:
        файлы метрик TPC-DS и потокового сценария в формате CSV, сценарий разбора
        и агрегации \texttt{parse\_writebench.py} (поддерживает вычисление как
        среднего, так и медианы).

    \item \textbf{\href{https://github.com/mai-diploma-kochkozharov/diploma-latex-template}{\texttt{diploma-latex-template}}}
        (ветка \texttt{vkr-kochkozharov}). Исходные тексты настоящей выпускной
        квалификационной работы: \LaTeX-шаблон оформления, исходные
        \texttt{.tex}-файлы по разделам, конфигурация сборки (\texttt{latexmk},
        \texttt{xelatex}), исходники диаграмм в формате draw.io и их PDF-сборки.
\end{itemize}

Доступ ко всем репозиториям открытый; клонирование и сборка артефактов возможны
без авторизации.

    \appendixsection{Фрагменты исходного кода}

В приложении приведены ключевые фрагменты разработанного модуля интеграции.
Рисунок~\ref{src:native} --- объявление нативных методов доступа к библиотеке Vortex;
рисунок~\ref{src:readbatch} --- метод чтения очередного пакета строк с экспортом в
раскладку Apache Arrow без копирования; рисунок~\ref{src:create} --- фрагмент
параметров создания таблицы near-real-time-сценария.

\begin{figure}
\begin{lstlisting}[language=Java]
/** Нативные методы доступа к файлам формата Vortex. */
public final class NativeFileMethods {
    static { NativeLoader.loadJni(); }
    private NativeFileMethods() {}

    public static native List<String> listVortexFiles(String uri, Map<String,String> options);
    public static native void   delete(String[] uris, Map<String,String> options);
    public static native long   open(String uri, Map<String,String> options);
    public static native long   rowCount(long pointer);
    public static native long   dtype(long pointer);
    public static native void   close(long pointer);
    public static native long   scan(long pointer, List<String> columns,
                                     byte[] predicateProto,
                                     long[] rowRange, long[] rowIndices);
}
\end{lstlisting}
\caption{Объявление нативных методов (\texttt{NativeFileMethods})}
\label{src:native}
\end{figure}

\begin{figure}
\begin{lstlisting}[language=Java]
@Override
public FileRecordIterator<InternalRow> readBatch() throws IOException {
    if (!arrayIterator.hasNext()) return null;
    releaseCurrentArray();
    Array array = arrayIterator.next();          // нативный пакет
    this.currentArray = array;
    VectorSchemaRoot vsr = array.exportToArrow(allocator, reuse); // Arrow C Data Interface, zero-copy
    this.reuse = vsr;
    Iterator<InternalRow> rows = arrowBatchReader.readBatch(vsr).iterator();
    return new FileRecordIterator<InternalRow>() {
        @Override public long returnedPosition() { return returnedPosition; }
        @Override public Path filePath()         { return filePath; }
        @Override public InternalRow next() {
            if (!rows.hasNext()) return null;
            returnedPosition = positionIterator.next();
            return rows.next();
        }
        @Override public void releaseBatch() { releaseCurrentArray(); }
    };
}
\end{lstlisting}
\caption{Чтение очередного пакета строк с экспортом в Apache Arrow (\texttt{VortexRecordsReader})}
\label{src:readbatch}
\end{figure}

\begin{figure}
\begin{lstlisting}[language=SQL]
CREATE TABLE bench_table ( /* 25 столбцов */, PRIMARY KEY (id) NOT ENFORCED) WITH (
  'file.format' = 'vortex',           -- варьируется: parquet | orc | vortex
  'bucket' = '...', 'bucket-key' = 'id',
  'sequence.field' = 'seq', 'rowkind.field' = 'op',
  'changelog-producer' = 'none', 'deletion-vectors.enabled' = 'false',
  'snapshot.num-retained.min' = '500', 'snapshot.num-retained.max' = '10000',
  'snapshot.time-retained' = '24h', 'snapshot.expire.execution-mode' = 'async',
  'manifest.merge-min-count' = '1000', 'consumer.expiration-time' = '24h',
  'write-buffer-size' = '1GB',
  'metadata.stats-mode' = 'full',     -- полные статистики (без усечения строк)
  'write.batch-size' = '65536'        -- крупный пакет: меньше пересечений границы JNI
);
\end{lstlisting}
\caption{Фрагмент параметров создания таблицы near-real-time-сценария}
\label{src:create}
\end{figure}

    \appendixsection{Параметры экспериментов}

В приложении приведены параметры окружения и запусков нагрузочного тестирования.

Кластер состоит из пяти узлов; на каждом --- два менеджера задач Apache Flink по 30
слотов (итого 10 менеджеров задач, 300 слотов); на одном узле дополнительно ---
менеджер заданий. Среда исполнения --- Java~17. Бэкенд хранения --- объектное
хранилище, совместимое с S3 (на базе Apache Ozone).
Развёртывание --- средствами Ansible.

В сценарии~1 (TPC-DS) масштаб набора --- \SI{100}{\giga\byte}; данные
загружены в Apache Paimon как таблицы только для добавления, разбиты на разделы
(типичный размер файла раздела --- порядка \SI{256}{\mega\byte}); параллелизм
запросов --- 300; на всех форматах выполнились \num{104} запроса (три прогона,
усреднение); сравнивались форматы Parquet, ORC, Vortex и Lance.

В сценарии~2 (потоковый near-real-time) параметры представительного запуска
приведены в таблице~\ref{tab:run-params}.

\begin{table}
    \caption{Параметры запуска потокового near-real-time-сценария}
    \label{tab:run-params}
    \begin{tabularx}{\textwidth}{|>{\raggedright\arraybackslash}p{0.32\textwidth}|X|}\hline
    \textbf{Параметр} & \textbf{Значение} \\ \hline
    Форматы & ORC, Parquet, Vortex \\ \hline
    Число бакетов & порядка сотни \\ \hline
    Фаза INIT & \num{50000000} строк, последовательные ключи, без ограничения скорости \\ \hline
    Фаза UPDATE & порядка \num{200000000} строк, случайные ключи, ограничение скорости $\approx 100\,000$~строк/с \\ \hline
    Интервал чекпойнтов Flink & 240~с (режим near-real-time) \\ \hline
    DURING-выборки & пакетные запросы во время записи с интервалом в несколько минут \\ \hline
    Число повторов цикла & не менее трёх \\ \hline
    Размер буфера записи & \SI{1}{\giga\byte} \\ \hline
    Размер пакета записи Vortex & \num{65536} строк \\ \hline
    Режим статистик & полные статистики по столбцам \\ \hline
    Политики хранения & расширенные (снимки 500--10000, время хранения 24~ч, асинхронная очистка) \\ \hline
    \end{tabularx}
\end{table}

Использованы следующие метрики. Для записи --- пропускная способность (строк в секунду), отдельно по
фазам INIT и UPDATE (с оговоркой о дедупликации в фазе UPDATE). Для чтения --- время
выполнения запроса в миллисекундах, отдельно по фазам DURING и FINAL, с усреднением
по выборкам и повторам. Дополнительно фиксировались число успешно завершённых
запросов под нагрузкой и наличие выбросов времени выполнения.

\end{document}
